% =====================================================================
% deephold_db — Handbuch
% Vollständige Dokumentation des Projekts in Lehrbuchform
% =====================================================================

\documentclass[11pt,a4paper,twoside,openright]{book}

% --- Sprache & Encoding ----------------------------------------------
\usepackage[utf8]{inputenc}
\usepackage[T1]{fontenc}
\usepackage[english]{babel}
\usepackage{csquotes}

% --- Layout & Typografie ---------------------------------------------
\usepackage[a4paper,margin=2.5cm,headheight=14pt]{geometry}
\usepackage{microtype}
\usepackage{parskip}
\usepackage{setspace}
\onehalfspacing

% --- Schriften --------------------------------------------------------
\usepackage{lmodern}
\usepackage[T1]{fontenc}
\usepackage{inconsolata}

% --- Grafik & Diagramme ----------------------------------------------
\usepackage{graphicx}
\usepackage{xcolor}
\usepackage{tikz}
\usetikzlibrary{shapes.geometric, arrows.meta, positioning, fit, backgrounds, calc, trees}
\usepackage{pgfplots}
\pgfplotsset{compat=1.18}

% --- Tabellen ---------------------------------------------------------
\usepackage{booktabs}
\usepackage{longtable}
\usepackage{array}
\usepackage{multirow}

% --- Math ---------------------------------------------------------------
\usepackage{amsmath}
\usepackage{amssymb}

% --- Code-Listings ----------------------------------------------------
\usepackage{listings}
\lstset{
  basicstyle=\ttfamily\small,
  keywordstyle=\color{blue!70!black}\bfseries,
  stringstyle=\color{red!60!black},
  commentstyle=\color{green!50!black}\itshape,
  numberstyle=\tiny\color{gray},
  numbers=left,
  numbersep=8pt,
  showstringspaces=false,
  breaklines=true,
  frame=single,
  framerule=0.4pt,
  rulecolor=\color{black!30},
  xleftmargin=2em,
  framexleftmargin=1.5em,
  backgroundcolor=\color{black!3},
  captionpos=b,
}
\lstdefinelanguage{yaml}{
  keywords={true,false,null,y,n},
  keywordstyle=\color{blue!70!black}\bfseries,
  sensitive=false,
  comment=[l]{\#},
  commentstyle=\color{green!50!black}\itshape,
  stringstyle=\color{red!60!black},
  morestring=[b]',
  morestring=[b]",
}
\lstdefinelanguage{sql}{
  keywords={SELECT,FROM,WHERE,JOIN,LEFT,INNER,OUTER,ON,AS,AND,OR,NOT,NULL,IS,IN,ORDER,BY,GROUP,HAVING,LIMIT,OFFSET,INSERT,INTO,VALUES,UPDATE,SET,DELETE,CREATE,TABLE,INDEX,UNIQUE,PRIMARY,KEY,FOREIGN,REFERENCES,DEFAULT,CHECK,CONSTRAINT,DROP,ALTER,ADD,COLUMN,UNION,ALL,INTERSECT,EXCEPT,CAST,AS,ASC,DESC,DISTINCT,BETWEEN,LIKE,ILIKE,SIMILAR,EXISTS,IN,ANY,SOME,CASE,WHEN,THEN,ELSE,END,COALESCE,NULLIF,GREATEST,LEAST,EXTRACT,DATE,PARTITION,OVER,ROW_NUMBER,RANK,DENSE_RANK,LAG,LEAD,FIRST_VALUE,LAST_VALUE,WITH,RECURSIVE},
  keywordstyle=\color{blue!70!black}\bfseries,
  sensitive=false,
  comment=[l]{--},
  commentstyle=\color{green!50!black}\itshape,
  stringstyle=\color{red!60!black},
  morestring=[b]',
  morestring=[b]",
}
\lstdefinelanguage{python}{
  keywords={False,None,True,and,as,assert,async,await,break,class,continue,def,del,elif,else,except,finally,for,from,global,if,import,in,is,lambda,nonlocal,not,or,pass,raise,return,try,while,with,yield,match,case},
  keywordstyle=\color{blue!70!black}\bfseries,
  sensitive=true,
  comment=[l]{\#},
  commentstyle=\color{green!50!black}\itshape,
  stringstyle=\color{red!60!black},
  morestring=[b]',
  morestring=[b]",
  morestring=[b]""",
}
\lstdefinelanguage{typescript}{
  keywords={abstract,as,async,await,boolean,break,case,catch,class,const,continue,debugger,declare,default,delete,do,else,enum,export,extends,false,finally,for,from,function,get,if,implements,import,in,instanceof,interface,is,keyof,let,module,namespace,never,new,null,number,object,package,private,protected,public,readonly,return,set,static,string,super,switch,this,throw,true,try,type,typeof,undefined,unique,unknown,var,void,while,with,yield,from,of},
  keywordstyle=\color{blue!70!black}\bfseries,
  sensitive=true,
  comment=[l]{//},
  commentstyle=\color{green!50!black}\itshape,
  morecomment=[s]{/*}{*/},
  stringstyle=\color{red!60!black},
  morestring=[b]',
  morestring=[b]",
  morestring=[b]`,
}

% --- Querverweise & Hyperlinks ---------------------------------------
\usepackage[hidelinks]{hyperref}
\hypersetup{
  colorlinks=true,
  linkcolor=blue!60!black,
  citecolor=blue!60!black,
  urlcolor=blue!60!black,
  pdftitle={deephold\_db -- Handbuch},
  pdfauthor={deephold\_db contributors},
}

% --- Beispiele und Boxen ---------------------------------------------
\usepackage{tcolorbox}
\tcbuselibrary{breakable, skins}

\newtcolorbox{merke}{enhanced, breakable, colback=yellow!8!white, colframe=yellow!50!black,
  fonttitle=\bfseries, title={Merke}}

\newtcolorbox{definition}{enhanced, breakable, colback=blue!5!white, colframe=blue!50!black,
  fonttitle=\bfseries, title={Definition}}

\newtcolorbox{warnung}{enhanced, breakable, colback=red!5!white, colframe=red!60!black,
  fonttitle=\bfseries, title={Achtung}}

\newtcolorbox{beispiel}{enhanced, breakable, colback=green!5!white, colframe=green!40!black,
  fonttitle=\bfseries, title={Beispiel}}

% --- Kopf- und Fußzeilen ---------------------------------------------
\usepackage{fancyhdr}
\pagestyle{fancy}
\fancyhf{}
\fancyhead[LE]{\nouppercase{\leftmark}}
\fancyhead[RO]{\nouppercase{\rightmark}}
\fancyfoot[LE,RO]{\thepage}
\renewcommand{\chaptermark}[1]{\markboth{Kapitel \thechapter\ -- #1}{}}
\renewcommand{\sectionmark}[1]{\markright{\thesection\ #1}}

% --- Inhaltsverzeichnis-Stil ----------------------------------------
\usepackage[titles]{tocloft}
\renewcommand{\cftchapfont}{\large\bfseries}
\renewcommand{\cftchappagefont}{\large}

% --- Literaturverzeichnis -------------------------------------------
\usepackage[backend=biber, style=alphabetic, sorting=ynt]{biblatex}
\addbibresource{literatur.bib}

% --- Index ------------------------------------------------------------
\usepackage{imakeidx}
\makeindex[intoc, columns=2, title=Index]

% --- Theorem-Umgebungen (für Begriffsdefinitionen) ------------------
\newtheorem{begriff}{Begriff}[chapter]

% --- Eigene Befehle --------------------------------------------------
\newcommand{\code}[1]{\texttt{\small #1}}
\newcommand{\filetype}[1]{\texttt{#1}}
\newcommand{\sql}[1]{\texttt{\small #1}}
\newcommand{\shell}[1]{\texttt{\small #1}}
\newcommand{\term}[1]{\textbf{\textit{#1}}}
\newcommand{\foreign}[1]{\textit{#1}}

% --- Titelseite -------------------------------------------------------
\title{
  \vspace{2cm}
  {\Huge\bfseries deephold\_db}\\[0.5cm]
  {\Large Ein Handbuch in Lehrbuchform}\\[0.5cm]
  {\large Vom Datenmodell bis zur Walk-Forward-Replikation\\
   einer quantitativen Forschungsarbeit}\\[2cm]
}
\author{deephold\_db contributors}
\date{\today}

\begin{document}

% --- Vorspann --------------------------------------------------------
\frontmatter
\maketitle
\thispagestyle{empty}

\vspace*{4cm}
\begin{center}
  {\Huge\bfseries deephold\_db}\\[0.5cm]
  {\Large Ein Handbuch in Lehrbuchform}\\[2cm]
  {\large Vom Datenmodell bis zur Walk-Forward-Replikation\\
  einer quantitativen Forschungsarbeit}\\[4cm]
  {\small Stand: \today}\\[4cm]
\end{center}

\vspace*{\fill}
\noindent\rule{\textwidth}{0.4pt}\\
{\footnotesize
Dieses Handbuch ist im Quellcode-Repository unter
\texttt{docs/handbuch/} abgelegt. Es ist als Lehrmaterial gedacht
und setzt Programmierkenntnisse voraus, aber keine
Finanzmarktexpertise.\\[0.5em]
Lizenz des Quellcodes: nicht-kommerziell, kein SLA (siehe \texttt{LICENSE} im Repo).
}

\cleardoublepage

% --- Vorwort --------------------------------------------------------
% =====================================================================
% Kapitel 0: Vorwort
% =====================================================================

\chapter{Vorwort}

\section{Wofür dieses Handbuch}

Dieses Handbuch dokumentiert das Projekt \code{deephold\_db} -- eine
Finanzmarktdatenbank für Forschung und privaten Gebrauch. Es ist nicht
nur eine technische Referenz, sondern als \emph{Lehrbuch} geschrieben:
Es erklärt Konzepte, rechtfertigt Designentscheidungen und führt den
Leser Schritt für Schritt durch die Architektur.

Das Handbuch richtet sich an drei Zielgruppen:

\begin{itemize}
  \item \textbf{Entwickler}, die das Projekt verstehen, erweitern oder
        debuggen wollen. Sie brauchen einen vollständigen Überblick über
        Architektur, Datenmodell, ETL-Pipeline, Tests und Deployment.
  \item \textbf{Forscher}, die quantitative Finanzanalysen auf Basis
        der gesammelten Daten durchführen wollen -- etwa Walk-Forward-Backtests
        wie in Kapitel~\ref{ch:analytics} beschrieben.
  \item \textbf{Studierende und Einsteiger}, die ein Gefühl dafür
        bekommen wollen, wie ein real-existierendes Datenprodukt
        entsteht -- vom Docker-Container bis zur Performance-Metrik.
\end{itemize}

\begin{merke}
Dieses Handbuch setzt Programmierkenntnisse voraus -- idealerweise
Python und Grundkonzepte relationaler Datenbanken -- aber \emph{keine}
Finanzmarktexpertise. Alle Fachbegriffe werden bei ihrer ersten
Verwendung in einer \textbf{Definitionsbox} (blau) erklärt.
\end{merke}

\section{Was \code{deephold\_db} nicht ist}

\begin{warnung}
Dieses Projekt ist ausdrücklich \textbf{kein} Bloomberg- oder
Refinitiv-Ersatz. Es deckt bewusst nur Retail- und Research-Bedürfnisse ab:

\begin{itemize}
  \item \textbf{Keine} echten Point-in-Time-Daten -- die adjusted-close-Reihen
        von Yahoo Finance sind \emph{mit aktuellem Wissensstand korrigiert}
        (look-ahead-Bias bei echten Backtests vor 2010 möglich).
  \item \textbf{Keine} CRSP-äquivalente Aktienverlauf-Datenbank -- kein
        Delisting-Tracking, keine Survivorship-Bias-Korrektur.
  \item \textbf{Keine} Intraday- oder Tick-Daten -- alles ist
        Tagesfrequenz.
  \item \textbf{Keine} echte Continuous-Future-Konstruktion -- Front-Month
        Futures werden einzeln betrachtet, kein Roll-Kalender.
  \item \textbf{Keine} Option-Chains, keine Bond-CUSIPs, keine
        OTC-Derivate.
\end{itemize}
\end{warnung}

Für institutionelle Anwendungen mit Echtzeitdaten, vollständiger
PIT-Historie und Compliance-Audit-Trail ist dieses Projekt nicht
geeignet. Es ist explizit als \emph{Selbstlern- und Forschungsplattform}
konzipiert.

\section{Wie dieses Handbuch zu lesen ist}

\begin{beispiel}
\textbf{Lese-Empfehlungen}
  \begin{itemize}
    \item \textbf{Entwickler} lesen das Handbuch linear von vorne nach
          hinten. Die Kapitel bauen aufeinander auf.
    \item \textbf{Forscher} springen direkt zu Kapitel~\ref{ch:analytics}
          (AEGIS-Replikation) und nutzen Glossar (Anhang~A) zum
          Nachschlagen.
    \item \textbf{Einsteiger} beginnen mit Kapitel~\ref{ch:grundlagen}
          (Grundbegriffe) und lesen dann selektiv die Kapitel, die sie
          interessieren. Das Glossar hilft beim Verständnis der
          Fachbegriffe.
  \end{itemize}
\end{beispiel}

\section{Aufbau}

Das Handbuch gliedert sich in zehn Hauptkapitel und vier Anhänge:

\begin{description}
  \item[Kapitel 1 -- Einführung] Was das Projekt tut, seine Ziele, sein
        Scope und seine Anwendungsfälle.
  \item[Kapitel 2 -- Grundlagen] Erklärung aller Fachbegriffe, die im
        restlichen Handbuch vorkommen -- von PostgreSQL bis zu
        CAGR.
  \item[Kapitel 3 -- Infrastruktur] Docker, Docker Compose, das Makefile,
        Python-Environment und die Bun-Runtime.
  \item[Kapitel 4 -- Datenmodell] Das ER-Diagramm, die 14~Tabellen und
        ihre Indizes.
  \item[Kapitel 5 -- Vendor-Adapter] Die drei Datenlieferanten FRED, ECB
        und Yahoo Finance.
  \item[Kapitel 6 -- ETL-Pipeline] Seeder und das \code{query\_db.py}-Skript.
  \item[Kapitel 7 -- Notebook-Generator] Das Plotly-Notebook und sein
        Build-Prozess.
  \item[Kapitel 8 -- OpenTUI-Explorer] Das Terminal-UI auf Basis von Bun
        und React 19.
  \item[Kapitel 9 -- Tests] pytest, \code{bun test} und die manuelle
        Test-Checklist.
  \item[Kapitel 10 -- Analytics \& AEGIS-Replikation] Die
        Walk-Forward-Backtest-Implementierung und ihre Ergebnisse
        im Vergleich zum Originalpaper.
  \item[Anhang A -- Glossar] Alphabetisches Verzeichnis aller
        Fachbegriffe.
  \item[Anhang B -- Befehlsreferenz] Alle Makefile-Targets und
        CLI-Befehle.
  \item[Anhang C -- Datenbankschema] Vollständiges DDL aller Tabellen.
  \item[Anhang D -- Literatur] Bibliographie der zitierten Werke.
\end{description}

\section{Danksagung}

Dieses Projekt wäre ohne die folgenden Open-Source-Bibliotheken und
Datenquellen nicht möglich:

\begin{itemize}
  \item \textbf{PostgreSQL 16} -- die Datenbank-Engine.
  \item \textbf{SQLAlchemy 2.x} -- das ORM.
  \item \textbf{Pandas, Polars, NumPy} -- Datenverarbeitung in Python.
  \item \textbf{httpx, tenacity, structlog} -- robustes HTTP und Logging.
  \item \textbf{Prefect 2.x} -- Workflow-Orchestrierung (vorbereitet).
  \item \textbf{Plotly} -- interaktive Plots im Browser.
  \item \textbf{Bun, React 19, OpenTUI} -- das TUI-Frontend.
  \item \textbf{FRED, ECB, Yahoo Finance} -- die drei Datenlieferanten.
  \item \textbf{AEGIS-Paper} -- der Anlass für die Backtest-Replikation.
\end{itemize}

Vielen Dank an alle Maintainer dieser Projekte.


% --- Inhaltsverzeichnis ----------------------------------------------
\tableofcontents

% --- Abbildungsverzeichnis (leer, da keine externen Bilder) ---------
\listoffigures
\addcontentsline{toc}{chapter}{\listfigurename}
\thispagestyle{empty}
\listoftables
\addcontentsline{toc}{chapter}{\listtablename}
\thispagestyle{empty}

% --- Hauptteil -------------------------------------------------------
\mainmatter

% =====================================================================
% Kapitel 1: Einführung und Projektüberblick
% =====================================================================

\chapter{Einführung und Projektüberblick}

\label{ch:einfuehrung}

\section{Das Problem}

Quantitative Finanzforschung lebt von Daten -- aber an genau diesen
Daten scheitern die meisten Hobby- und Kleinprojekte:

\begin{itemize}
  \item \textbf{Datenquellen sind fragmentiert}. Makro-Daten kommen vom
        FED, Aktien-Daten von Yahoo, Anleihe-Daten vom Treasury, FX-Daten
        von der EZB, Rohstoff-Daten von der COMEX. Jede Quelle hat ein
        anderes Format, eine andere API, andere Lizenzbedingungen.
  \item \textbf{Historie ist teuer}. Bloomberg-Terminals kosten
        20.000\,{}/Jahr, Refinitiv vergleichbar. Für
        Hobby-Forscher unbezahlbar.
  \item \textbf{Open Data ist da, aber unhandlich}. FRED, EZB SDMX,
        Yahoo Finance -- alle frei verfügbar, aber jede mit eigener
        API-Linie, eigenen Rate-Limits, eigenem Datenformat.
  \item \textbf{Reproduzierbarkeit ist schwer}. Wenn der Code zur
        Datensammlung nicht versioniert und mitprotokolliert ist,
        weiß man sechs Monate später nicht mehr, welcher Pull
        welche Daten geliefert hat.
  \item \textbf{Backtesting braucht saubere PIT-Daten}. Eine
        Aktienkurs-Zeitreihe vom 1.~Januar 2020 muss zu jedem
        späteren Zeitpunkt exakt so aussehen, wie sie tatsächlich
        an diesem Tag bekannt war -- inklusive aller nachträglichen
        Splits, Dividenden und Restatements.
\end{itemize}

\section{Die Lösung: \code{deephold\_db}}

\code{deephold\_db} ist eine selbst-gehostete, lokal laufende
Finanzmarktdatenbank, die alle genannten Probleme auf eine pragmatische
Art adressiert:

\begin{definition}
\textbf{Retail-Research-Datenbank}
Eine Datenbank, die \emph{für den privaten und Forschungsgebrauch}
konzipiert ist. Sie ist nicht für Hochfrequenzhandel, regulatorische
Compliance oder institutionelle Vermögensverwaltung gedacht. Sie
priorisiert \textbf{einfache Bedienung} und \textbf{kostengünstigen
Betrieb} über \textbf{Latenz} und \textbf{vollständige PIT-Treue}.
\end{definition}

Die Kernidee: \emph{ein Befehl}, \emph{eine Datenbank}, \emph{mehrere
Datenquellen}, \emph{mehrere Frontends}. Konkret bedeutet das:

\begin{enumerate}
  \item \textbf{Ein-Befehl-Setup}. \code{make init} startet Docker,
        installiert Python-Dependencies, legt das Datenbankschema an
        und seedet die ersten Daten. Alles in unter 10~Minuten.
  \item \textbf{Mehrere Datenquellen, einheitliches Schema}. FRED
        (FED-Wirtschaftsdaten), EZB SDMX (EUR-Geldmarkt, EUR-Anleihen,
        EUR-FX) und Yahoo Finance (US-Aktien, ETFs, Indizes) werden über
        Vendor-Adapter in dasselbe Schema gemappt.
  \item \textbf{Mehrere Frontends}.
        \begin{itemize}
          \item \textbf{Jupyter-Notebook} (\code{build\_notebook.py})
                mit Plotly für interaktive Charts.
          \item \textbf{OpenTUI-Explorer} (Bun + React 19) für
                schnelles Durchsuchen der Datenbank im Terminal.
          \item \textbf{SQL direkt} via psql für Ad-hoc-Analysen.
          \item \textbf{Python-API} (SQLAlchemy ORM) für reproduzierbare
                Research-Skripte.
        \end{itemize}
  \item \textbf{Reproduzierbarkeit}. Jede Zeile in der Datenbank trägt
        einen Zeitstempel (\code{ingested\_at}) und eine
        \code{vendor\_id}. Damit lässt sich auch Monate später
        rekonstruieren, woher ein bestimmter Wert stammt.
  \item \textbf{Walk-Forward-Backtest-Framework}. Eine
        vereinfachte Version des AEGIS-Algorithmus aus dem
        quantitativen Finanzpaper von Chakraborty und Singh
        (2026) demonstriert, wie die gesammelten Daten für
        systematische Strategieentwicklung genutzt werden können.
\end{enumerate}

\section{Anwendungsfälle}

\begin{beispiel}
\textbf{Typische Anwendungsfälle}
  \begin{itemize}
    \item \textbf{Portfolio-Analyse}: ``Wie hat sich mein
          60/40-Aktien/Anleihen-Portfolio in den letzten 20~Jahren
          entwickelt?''
    \item \textbf{Makro-Research}: ``Wie korreliert die
          US-Arbeitslosenquote mit S\&P-500-Renditen in den
          letzten 10~Jahren?''
    \item \textbf{Faktor-Analyse}: ``Lohnt sich Value-Investing
          in Small-Caps noch?''
    \item \textbf{Backtest-Vorbereitung}: ``Lade mir 20~Jahre
          S\&P~500- und NASDAQ-Daten, damit ich meinen
          Mean-Reversion-Filter darauf testen kann.''
    \item \textbf{Bildungsprojekte}: ``Ich will meinen Studierenden
          zeigen, wie ein professionelles Datenprodukt
          aussieht.''
  \end{itemize}
\end{beispiel}

\section{Nicht-Ziele}

\begin{warnung}
\rmfamily\itshape
Folgende Funktionen sind \emph{bewusst nicht} Teil dieses Projekts:

\begin{itemize}
  \item Echtzeitdaten oder Intraday-Ticks.
  \item Echtes Point-in-Time mit allen Splits/Dividenden zum
        jeweiligen historischen Zeitpunkt (nur adjusted-close,
        keine nachträglichen Restatements).
  \item Bond-CUSIP-Level-Daten (nur aggregierte Anleihe-ETFs
        oder Makro-Yields).
  \item Multi-Vendor-Reconciliation (Preisvergleiche zwischen
        Datenquellen).
  \item Historische Index-Mitgliedschaften (``Wann wurde AAPL
        in den S\&P~500 aufgenommen?'').
  \item Historische Credit-Ratings.
\end{itemize}
\end{warnung}

\section{Architektur auf einen Blick}

\begin{figure}[ht]
\centering
\begin{tikzpicture}[
  node distance=0.8cm and 1.2cm,
  box/.style={rectangle, draw, thick, rounded corners, minimum width=2.8cm, minimum height=0.8cm, align=center, font=\small},
  layer/.style={rectangle, draw, very thick, rounded corners, fill=#1!10, inner sep=0.4cm},
  arrow/.style={-Latex, thick},
]
  \node[box, fill=blue!20] (python) {Python 3.11\\(SQLAlchemy, Polars)};
  \node[box, fill=green!20, right=of python] (db) {PostgreSQL 16\\(14 Tabellen)};
  \node[box, fill=yellow!20, right=of db] (bun) {Bun + React 19\\(OpenTUI)};

  \node[box, fill=red!20, below=of python] (vendors) {Vendor-Adapter\\(FRED, ECB, Yahoo)};
  \node[box, fill=cyan!20, below=of db] (etl) {ETL-Pipeline\\(Seeder, query\_db.py)};
  \node[box, fill=orange!20, below=of bun] (notebook) {Plotly-Notebook\\(build\_notebook.py)};

  \draw[arrow] (vendors) -- node[above, font=\scriptsize] {insert} (db);
  \draw[arrow] (etl) -- node[above, font=\scriptsize] {read/write} (db);
  \draw[arrow] (python) -- node[above, font=\scriptsize] {SQLAlchemy} (db);
  \draw[arrow] (bun) -- node[above, font=\scriptsize] {node-postgres} (db);
  \draw[arrow] (notebook) -- node[above, font=\scriptsize] {read-only} (db);
\end{tikzpicture}
\caption{High-Level-Architektur von \code{deephold\_db}: Python-Backend
  füllt die DB über Vendor-Adapter, das TUI und das Notebook lesen
  direkt aus der DB. Die einzige Kopplung zwischen den Layern ist
  die Datenbank selbst.}
\label{fig:arch-overview}
\end{figure}

Abbildung~\ref{fig:arch-overview} zeigt die drei Schichten des
Projekts:

\begin{enumerate}
  \item \textbf{Python-Backend} (links): Enthält die Vendor-Adapter
        (FRED, ECB, Yahoo), die ETL-Pipeline (Seeder, \code{query\_db.py})
        und die Analytics-Module (Returns, Metriken, VAM-Backtest).
        Spricht via SQLAlchemy mit der DB.
  \item \textbf{PostgreSQL-Datenbank} (Mitte): Single Source of Truth.
        14~Tabellen -- \code{venues}, \code{vendors}, \code{instruments},
        \code{instrument\_identifiers}, \code{trading\_calendars},
        \code{prices\_raw}, \code{prices\_daily}, \code{corporate\_actions},
        \code{fx\_rates\_daily}, \code{bond\_yields}, \code{macro\_series},
        \code{macro\_observations}, \code{dq\_runs}, \code{dq\_findings}.
  \item \textbf{TypeScript-TUI} (rechts): OpenTUI-Explorer in Bun.
        Liest die DB via \code{node-postgres}, kein Python im Hot-Path.
        Maximale Performance für den Endbenutzer.
\end{enumerate}

\section{Was Sie als Nächstes erwartet}

Das nächste Kapitel~--\ref{ch:grundlagen}--~bringt Sie auf den
nötigen Wissensstand, falls Sie mit einem der verwendeten Konzepte
noch nicht vertraut sind. Falls Sie bereits Erfahrung mit PostgreSQL,
Docker, Python und TypeScript haben, können Sie direkt mit
Kapitel~\ref{ch:infrastruktur} fortfahren.

% =====================================================================
% Kapitel 2: Grundlagen
% =====================================================================

\chapter{Grundlagen}

\label{ch:grundlagen}

Dieses Kapitel erklärt alle Fachbegriffe, die in den Folgekapiteln
verwendet werden. Wer mit allen Konzepten vertraut ist, kann es
überspringen und bei Unklarheiten als Nachschlagewerk nutzen.

\section{Datenbank-Grundlagen}

\begin{definition}
\textbf{Relationale Datenbank}
Eine \emph{relationale Datenbank} organisiert Daten in
\emph{Tabellen} (auch \emph{Relationen} genannt) mit Zeilen (Datensätze)
und Spalten (Attribute). Zwischen den Tabellen bestehen
\emph{Beziehungen} über \emph{Fremdschlüssel} (\foreign{foreign key}).
Die meisten modernen Datenbanken -- einschließlich PostgreSQL --
sprechen die standardisierte Abfragesprache SQL.
\end{definition}

\begin{definition}
\textbf{SQL (Structured Query Language)}
SQL ist die standardisierte Abfragesprache für relationale Datenbanken.
SQL unterteilt sich in Untergruppen:

\begin{itemize}
  \item \textbf{DDL} (Data Definition Language): \code{CREATE TABLE},
        \code{ALTER TABLE}, \code{DROP TABLE}.
  \item \textbf{DML} (Data Manipulation Language): \code{SELECT},
        \code{INSERT}, \code{UPDATE}, \code{DELETE}.
  \item \textbf{DCL} (Data Control Language): \code{GRANT},
        \code{REVOKE}.
\end{itemize}
\end{definition}

\begin{definition}
\textbf{PostgreSQL}
PostgreSQL -- oft kurz \emph{Postgres} -- ist ein
\emph{objekt-relationales} Datenbankmanagementsystem (ORDBMS). Es ist
quelloffen, kostenlos, seit über 30~Jahren in Entwicklung und gilt als
eine der stabilsten und funktionsreichsten Open-Source-Datenbanken.
PostgreSQL unterstützt JSON/JSONB, Volltextsuche, Window Functions und
vieles mehr.
\end{definition}

\begin{beispiel}
\textbf{PostgreSQL-Verbindung in Python}
\begin{lstlisting}[language=python]
import psycopg

conn = psycopg.connect(
    host="localhost",
    port=5432,
    user="deephold",
    password="deephold",
    dbname="deephold",
)
with conn.cursor() as cur:
    cur.execute("SELECT COUNT(*) FROM macro_observations")
    print(cur.fetchone())
\end{lstlisting}
\end{beispiel}

\begin{definition}
\textbf{Index}
Ein \emph{Index} ist eine zusätzliche Datenstruktur, die das schnelle
Suchen in einer Tabelle ermöglicht -- analog zum Stichwortverzeichnis
eines Buchs. Ein Index auf Spalte~\code{X} beschleunigt Abfragen
mit \code{WHERE X = ...}, verlangsamt aber \code{INSERT}-Operationen
und kostet zusätzlichen Speicherplatz. Daher werden nur Spalten
indiziert, die häufig in \code{WHERE}-Klauseln vorkommen.
\end{definition}

\begin{beispiel}
\textbf{Index-Definition}
\begin{lstlisting}[language=sql]
CREATE INDEX ix_prices_daily_date
  ON prices_daily (date);

CREATE UNIQUE INDEX ix_instrument_identifiers_scheme_value
  ON instrument_identifiers (scheme, value);
\end{lstlisting}
\end{beispiel}

\begin{definition}
\textbf{JSONB}
JSONB ist die binäre Variante des JSON-Datentyps in PostgreSQL. Werte
werden nicht als Text, sondern als effizient geparste interne Struktur
gespeichert. Das ermöglicht Indexierung auf JSON-Feldern und
schnelle Pfad-Queries (\code{data-\texttt{>}'field'}).
\end{definition}

\section{Python-Grundlagen}

\begin{definition}
\textbf{Python}
Python ist eine interpretierte, dynamisch typisierte
Hochsprachen-Programmiersprache. Sie wird im Data-Science-Umfeld
besonders wegen ihrer einfachen Syntax, ihrer umfangreichen
Standardbibliothek und ihres dichten Ökosystems an
Dritt-Bibliotheken geschätzt. Das Projekt verwendet Python in
Version~3.11 oder neuer.
\end{definition}

\begin{definition}
\textbf{Poetry}
Poetry ist ein modernes Tool für \emph{Dependency-Management} in
Python-Projekten. Es verwaltet die \code{pyproject.toml}-Datei, legt
virtuelle Environments automatisch an und löst transitive
Abhängigkeiten auf. In diesem Projekt wird es durch
\code{pip} und ein venv-Verzeichnis ersetzt (siehe Kapitel~\ref{ch:infrastruktur}),
weil Poetry nicht überall verfügbar ist.
\end{definition}

\begin{definition}
\textbf{Pandas}
Pandas ist die Standard-Bibliothek für tabellarische Datenverarbeitung
in Python. Die zentrale Datenstruktur ist das \emph{DataFrame} -- eine
beschriftete, zweidimensionale Tabelle. Pandas wird für Daten-
Cleaning, Resampling, Grouping und vieles mehr verwendet.
\end{definition}

\begin{definition}
\textbf{Polars}
Polars ist eine neuere, schnellere Alternative zu Pandas, implementiert
in Rust. Polars nutzt \emph{Lazy Evaluation} und \emph{Columnar
Storage}, was es auf großen Datensätzen drastisch schneller macht.
Die API ähnelt Pandas, ist aber stärker typisiert.
\end{definition}

\begin{beispiel}
\textbf{Polars vs.\ Pandas -- ein einfaches Beispiel}
\begin{lstlisting}[language=python]
import polars as pl
import pandas as pd

# Beide erzeugen das gleiche Ergebnis:
df_pl = pl.DataFrame({"a": [1, 2, 3], "b": [4.0, 5.0, 6.0]})
df_pd = pd.DataFrame({"a": [1, 2, 3], "b": [4.0, 5.0, 6.0]})

# Polars:
mean_b = df_pl["b"].mean()  # Lazy, gibt eine Expression zurück
print(df_pl.select(pl.col("b").mean()))  # 5.0

# Pandas:
mean_b = df_pd["b"].mean()  # Eager, gibt einen float zurück
print(mean_b)  # 5.0
\end{lstlisting}
\end{beispiel}

\begin{definition}
\textbf{SQLAlchemy}
SQLAlchemy ist ein \emph{ORM-Toolkit} (Object-Relational Mapper) für
Python. Es erlaubt, Datenbank-Tabellen als Python-Klassen zu
modellieren und Abfragen in Python-Syntax zu schreiben, ohne rohes
SQL zu schreiben. SQLAlchemy unterstützt zwei Stile:

\begin{itemize}
  \item \textbf{Core}: niedrigschwellige SQL-Bausteine
        (\code{select()}, \code{insert()} etc.).
  \item \textbf{ORM}: hochschwellige, objekt-orientierte
        Zugriffsschicht mit Klassen für jede Tabelle.
\end{itemize}
In diesem Projekt wird die ORM-Schicht verwendet.
\end{definition}

\begin{beispiel}
\textbf{SQLAlchemy-Definition einer Tabelle}
\begin{lstlisting}[language=python]
from sqlalchemy.orm import DeclarativeBase, Mapped, mapped_column

class Base(DeclarativeBase):
    pass

class Instrument(Base):
    __tablename__ = "instruments"
    instrument_id: Mapped[int] = mapped_column(primary_key=True)
    asset_class: Mapped[str] = mapped_column()
    name: Mapped[str] = mapped_column()
    # ...
\end{lstlisting}
\end{beispiel}

\begin{definition}
\textbf{Alembic}
Alembic ist das offizielle Migrationswerkzeug für SQLAlchemy. Es
verwaltet \emph{Schema-Migrationen} -- das Hinzufügen, Ändern oder
Löschen von Tabellen und Spalten, ohne Daten zu verlieren. Jede
Migration ist eine Python-Datei mit zwei Funktionen: \code{upgrade()}
und \code{downgrade()}.
\end{definition}

\begin{beispiel}
\textbf{Alembic-Migration -- sehr vereinfacht}
\begin{lstlisting}[language=python]
def upgrade() -> None:
    op.create_table(
        "macro_series",
        sa.Column("series_id", sa.Text, primary_key=True),
        sa.Column("name", sa.Text),
        # ...
    )

def downgrade() -> None:
    op.drop_table("macro_series")
\end{lstlisting}
\end{beispiel}

\section{Docker und Container}

\begin{definition}
\textbf{Container}
Ein \emph{Container} ist eine isolierte Laufzeit-Umgebung für eine
Anwendung. Anders als bei virtuellen Maschinen wird nicht der
gesamte Rechner virtualisiert, sondern nur die
\emph{Anwendungs-Layer} (Binaries, Libraries, Konfiguration).
Container starten in Sekunden, nicht in Minuten, und sind sehr
ressourcenschonend.
\end{definition}

\begin{definition}
\textbf{Docker}
Docker ist die populärste Container-Engine. Sie packt eine Anwendung
in ein \emph{Image} (eine Art Bauplan) und startet daraus
\emph{Container}. Ein Image wird durch ein \code{Dockerfile}
beschrieben, ein Container ist eine laufende Instanz eines Images.
\end{definition}

\begin{definition}
\textbf{Docker Compose}
Docker Compose ist ein Werkzeug zur Definition und zum Start
\emph{mehrerer} Container als zusammengehörige Anwendung. Die
Konfiguration erfolgt in einer \code{docker-compose.yml}-Datei.
Mit \code{docker compose up -d} startet man alle Services im
Hintergrund; mit \code{docker compose down} stoppt man sie wieder.
\end{definition}

\begin{beispiel}
\textbf{docker-compose.yml -- Auszug}
\begin{lstlisting}[language=yaml]
services:
  postgres:
    image: postgres:16-alpine
    environment:
      POSTGRES_USER: deephold
      POSTGRES_PASSWORD: deephold
      POSTGRES_DB: deephold
    ports:
      - "5432:5432"
    healthcheck:
      test: ["CMD-SHELL", "pg_isready -U deephold"]
      interval: 5s
\end{lstlisting}
\end{beispiel}

\begin{definition}
\textbf{Docker Volume}
Ein \emph{Volume} ist ein persistenter Speicherbereich, der vom
Container unabhängig existiert. Wird ein Container gelöscht, bleiben
Daten im Volume erhalten. \code{deephold\_db} definiert ein Volume
\code{postgres\_data} für die Datenbank, damit die Daten
Container-Updates überleben.
\end{definition}

\section{Networking-Grundlagen}

\begin{definition}
\textbf{HTTP (Hypertext Transfer Protocol)}
HTTP ist das Standard-Protokoll für die Kommunikation zwischen
Web-Clients und Web-Servern. Datenquellen wie FRED und Yahoo Finance
bieten ihre Daten über \emph{RESTful HTTP}-APIs an. Ein typischer
Aufruf ist \code{GET
https://api.example.com/data?param=value}, der mit einem
JSON-Array oder XML-Dokument beantwortet wird.
\end{definition}

\begin{definition}
\textbf{Rate Limiting}
Die meisten öffentlichen APIs begrenzen, wie viele Anfragen ein
einzelner Client pro Zeiteinheit stellen darf -- das sogenannte
\emph{Rate Limiting}. Wird das Limit überschritten, antwortet der
Server mit HTTP 429 (Too Many Requests). Das Projekt implementiert
ein \emph{Token-Bucket}-Verfahren, bei dem jeder Vendor ein eigenes
Bucket mit konfigurierbarer Rate und Burst-Größe hat.
\end{definition}

\begin{beispiel}
\textbf{Token-Bucket in Python}
\begin{lstlisting}[language=python]
import time

class TokenBucket:
    def __init__(self, rate_per_min: int, burst: int):
        self.rate = rate_per_min / 60.0   # tokens per second
        self.burst = burst
        self.tokens = float(burst)
        self.last = time.monotonic()

    def take(self, n: int = 1) -> None:
        now = time.monotonic()
        elapsed = now - self.last
        self.last = now
        self.tokens = min(self.burst, self.tokens + elapsed * self.rate)
        if self.tokens < n:
            # Blockierend warten
            time.sleep((n - self.tokens) / self.rate)
            self.tokens = 0
        else:
            self.tokens -= n
\end{lstlisting}
\end{beispiel}

\section{TypeScript- und TUI-Grundlagen}

\begin{definition}
\textbf{TypeScript}
TypeScript ist eine Obermenge von JavaScript, die statische Typen
hinzufügt. TypeScript-Code wird zu JavaScript compiliert und läuft
dann in jeder JavaScript-Engine (z.\,B.\ Node.js oder Bun).
\end{definition}

\begin{definition}
\textbf{Bun}
Bun ist eine moderne JavaScript-/TypeScript-Runtime, ähnlich wie
Node.js, aber deutlich schneller beim Start und bei der Installation
von Dependencies. Bun enthält einen eingebauten
\emph{Paketmanager}, einen \emph{Test-Runner} und einen
\emph{Bundler}. Bun ist die Standard-Runtime für das in diesem
Projekt verwendete OpenTUI-Framework.
\end{definition}

\begin{definition}
\textbf{Terminal User Interface (TUI)}
Eine \emph{TUI} ist ein textbasiertes Interface, das im Terminal
läuft -- im Gegensatz zu einem GUI (grafisches User Interface).
Im Gegensatz zu einem Kommandozeilen-Tool, das nur Text ein- und
ausgibt, bietet eine TUI ein \emph{interaktives} Layout mit
\emph{Widgets} wie Listen, Buttons, Formularfeldern und Charts.
Beispiele: vim, htop, lazygit.
\end{definition}

\begin{definition}
\textbf{OpenTUI}
OpenTUI ist ein von Anomaly (https://anoma.ly) entwickeltes
TUI-Framework. Es besteht aus einem in Zig geschriebenen
\emph{Core} und TypeScript-Bindings mit React- und Solid.js-Support.
Der Kern nutzt Yoga (Facebooks Flexbox-Engine), um Layouts zu
berechnen -- genau wie im Web -- und rendert das Ergebnis in einen
Terminal-Buffer.
\end{definition}

\begin{definition}
\textbf{Yoga (Flexbox-Layout)}
Yoga ist Facebooks plattformübergreifender Layout-Algorithmus, der
\foreign{flexbox} -- das CSS-Layout-System -- in native Bibliotheken
portiert. OpenTUI nutzt Yoga, um komplexe Layouts (Sidebars,
Header, Footer, mehrspaltige Inhalte) pixelgenau im Terminal zu
positionieren.
\end{definition}

\section{Finanzmarkt-Grundlagen}

\begin{definition}
\textbf{Adjusted Close}
Der \emph{adjusted Close} ist der Schlusskurs eines Wertpapiers,
\emph{nachbereinigt} um alle Kapitalmaßnahmen, die zwischen
Kaufzeitpunkt und aktuellem Zeitpunkt stattgefunden haben -- vor
allem Splits und Dividenden. Wenn eine Aktie am Tag~X zu 100\,{}
schließt und am Folgetag einen 2:1-Split macht, dann steht der
adjusted Close für Tag~X nicht bei 100, sondern bei 50\,{}.
Backtests, die adjusted close verwenden, sind -- in der Retail-Praxis
-- korrekt für historische Investoren.
\end{definition}

\begin{definition}
\textbf{Log-Return vs.\ Simple Return}
Es gibt zwei gängige Arten, Renditen zu berechnen:

\begin{itemize}
  \item \textbf{Simple Return}: $R_t = \frac{P_t - P_{t-1}}{P_{t-1}}$
  \item \textbf{Log-Return}: $r_t = \ln\frac{P_t}{P_{t-1}}$
\end{itemize}

Log-Returns sind \emph{zeitadditiv} -- der Log-Return über zwei
Tage ist die Summe der täglichen Log-Returns. Dies vereinfacht
viele Berechnungen. Bei kleinen Returns (~$|r_t| < 5\%$) sind die
beide Werte praktisch identisch.
\end{definition}

\begin{definition}
\textbf{Volatilität (Vol)}
Die \emph{Volatilität} ist die annualisierte Standardabweichung der
Renditen. Sie ist das Standardmaß für das Risiko eines Wertpapiers
oder Portfolios. Übliche Symbole: $\sigma$ (Sigma) oder $\text{vol}$.
\end{definition}

\begin{beispiel}
\textbf{Volatilität berechnen}
\begin{lstlisting}[language=python]
import math
import statistics

# Tägliche Log-Returns einer Aktie
daily_returns = [0.01, -0.02, 0.005, 0.015, -0.01, 0.0, 0.02]

# Standardabweichung der täglichen Returns
daily_vol = statistics.stdev(daily_returns)  # ≈ 0.0124

# Annualisiert (252 Handelstage)
annual_vol = daily_vol * math.sqrt(252)     # ≈ 0.196 = 19.6%
\end{lstlisting}
\end{beispiel}

\begin{definition}
\textbf{CAGR (Compound Annual Growth Rate)}
Die \emph{CAGR} ist die durchschnittliche jährliche Wachstumsrate
eines Investments über einen bestimmten Zeitraum, unter Annahme
eines geometrischen (zinseszinsartigen) Wachstums:

\[
  \text{CAGR} = \left(\frac{V_\text{end}}{V_\text{start}}\right)^{1/n} - 1
\]

wobei $n$ die Anzahl der Jahre ist. Ein Investment, das von 100\,{}
auf 200\,{} in 5~Jahren wächst, hat eine CAGR von
$\sqrt[5]{200/100} - 1 \approx 14{,}87\%$.
\end{definition}

\begin{definition}
\textbf{Sharpe Ratio}
Die \emph{Sharpe Ratio} misst die \emph{renditebereinigte Performance}
eines Investments -- die Überrendite pro Einheit Gesamtrisiko:

\[
  \text{Sharpe} = \frac{\bar{R}_p - R_f}{\sigma_p}
\]

wobei $\bar{R}_p$ die annualisierte Portfoliorendite, $R_f$ der
risikofreie Zins (oft 4\,\% p.\,a.) und $\sigma_p$ die annualisierte
Volatilität ist. Höher ist besser; eine Sharpe~>~1 gilt als gut,
Sharpe~>~2 als exzellent.
\end{definition}

\begin{definition}
\textbf{Sortino Ratio}
Eine Variante der Sharpe Ratio, die nur die
\emph{Abwärtsvolatilität} (negative Renditen) als Risikomaß
verwendet. Sie ist aussagekräftiger für asymmetrische Strategien, die
gerne hohe positive Volatilität (große Gewinne) in Kauf nehmen, aber
negative Volatilität (Verluste) minimieren wollen.
\end{definition}

\begin{definition}
\textbf{Maximum Drawdown (MaxDD)}
Der \emph{Maximum Drawdown} ist der größte prozentuale Verlust vom
einem vorherigen Höchststand bis zu einem nachfolgenden Tiefststand:

\[
  \text{MaxDD} = \min_t \frac{V_t - V_\text{peak}}{V_\text{peak}}
\]

Ein MaxDD von $-30\,\%$ bedeutet, dass das Portfolio irgendwann
30\,\% unter seinem vorherigen Höchststand lag. Dies ist das
wichtigste \emph{Tail-Risk}-Maß.
\end{definition}

\begin{definition}
\textbf{Walk-Forward-Backtest}
Ein \emph{Walk-Forward-Backtest} ist die realistischste Form der
Backtest-Validierung. Anders als bei einem klassischen
\emph{In-Sample}-Backtest wird das Modell nicht auf denselben Daten
trainiert und evaluiert, auf denen es entwickelt wurde. Stattdessen:

\begin{enumerate}
  \item Trainiere auf $[t - L, t]$ (z.\,B.\ 12~Monate)
  \item Evaluiere auf $[t, t+1]$ (z.\,B.\ 1~Monat)
  \item Wandere ein Fenster weiter, wiederhole.
\end{enumerate}

So bekommt das Modell zu jedem Zeitpunkt nur Daten, die es auch
\emph{tatsächlich} gehabt hätte -- kein \emph{Look-Ahead-Bias}.
\end{definition}

\begin{definition}
\textbf{Look-Ahead-Bias}
\emph{Look-Ahead-Bias} entsteht, wenn ein Backtest Informationen
verwendet, die zum Entscheidungszeitpunkt in der Realität noch nicht
verfügbar gewesen wären -- zum Beispiel ein Split, der erst nach dem
Backtest-Datum in den Preisdaten nachträglich eingerechnet wurde.
Walk-Forward-Tests eliminieren Look-Ahead-Bias vollständig.
\end{definition}

\begin{definition}
\textbf{Survivorship-Bias}
\emph{Survivorship-Bias} ist die Tendenz, nur die
\emph{überlebenden} Unternehmen in einem Index zu betrachten. Wenn
Sie z.\,B.\ die durchschnittliche Rendite der heutigen S\&P-500-Mitglieder
in der Vergangenheit berechnen, ignorieren Sie all die Unternehmen,
die zwischenzeitlich pleitegingen oder delistet wurden. Das
überschätzt die wahre historische Performance.
\end{definition}

\begin{definition}
\textbf{Momentum}
\emph{Momentum} ist die Tendenz von Aktien, die in der Vergangenheit
gut gelaufen sind, dies auch in der nahen Zukunft zu tun -- und
umgekehrt. Die klassische Momentum-Strategie (Jegadeesh \&
Titman, 1993) kauft die Top-10\,\% der zurückliegenden 12~Monate
und leert die unteren 10\,\%. Momentum ist eines der am besten
dokumentierten Anomalien in der Finanzmarktforschung.
\end{definition}

\begin{definition}
\textbf{VAM (Volatility-Adjusted Momentum)}
Die \emph{VAM} (manchmal auch \emph{Sharpe-Momentum} genannt)
normalisiert die historische Rendite durch die historische
Volatilität:

\[
  \text{VAM}_i = \frac{R_i}{\sigma_i}
\]

Eine Aktie mit 50\,\% Jahresrendite und 30\,\% Volatilität hat
einen VAM von $1{,}67$. Eine andere mit 50\,\% Rendite aber
60\,\% Volatilität hat einen VAM von $0{,}83$ -- sie ist bei
gleicher Rendite deutlich riskanter. VAM hilft, den
\emph{Winner's Curse} zu vermeiden: Aktien, die nur wegen hoher
Beta gut laufen, werden aussortiert.
\end{definition}

\begin{definition}
\textbf{Skip-Month}
Beim Walk-Forward-Backtest ist es üblich, den letzten Monat vor
dem Re-Balancing-Datum \emph{auszulassen} (das sogenannte
\emph{Skip-Month}). Grund: Die jüngsten 21~Tage enthalten oft
Mean-Reversion-Effekte und Mikrostruktur-Rauschen, die die
Signale verzerren würden. Das AEGIS-Paper und dieses Projekt
verwenden ein 21-Tage-Skip (Default).
\end{definition}

\begin{definition}
\textbf{Friction / Transaction Costs}
\emph{Friction} (auch \emph{Transaction Costs}) sind die
\emph{Transaktionskosten} eines Backtests -- typischerweise als
Bps (Basispunkte, 1~Bp = 0{,}01\,\%) pro umgesetztem Dollar
ausgedrückt. Das AEGIS-Paper verwendet 10~Bp pro Handel, was
etwa 1~Bp pro Monat-Portfolio-Turnover entspricht.
\end{definition}

\section{Was kommt als Nächstes?}

Mit diesen Grundlagen bewaffnet, können Sie sich auf die
Infrastruktur (Kapitel~\ref{ch:infrastruktur}) stürzen. Dort
sehen wir, wie alle diese Komponenten in einem lauffähigen
System zusammenspielen.

% =====================================================================
% Kapitel 3: Infrastruktur
% =====================================================================

\chapter{Infrastruktur}

\label{ch:infrastruktur}

Dieses Kapitel zeigt, wie alle Komponenten -- Datenbank, Python,
TypeScript, Tools -- auf einer Maschine zusammenstarten und
miteinander kommunizieren. Es ist die Grundlage, auf der alle
folgenden Kapitel aufbauen.

\section{Überblick: Was läuft wo?}

Abbildung~\ref{fig:infra-overview} zeigt, welche Prozesse auf der
Entwicklungsmaschine laufen und wie sie miteinander sprechen.

\begin{figure}[ht]
\centering
\begin{tikzpicture}[
  node distance=1cm and 1.5cm,
  box/.style={rectangle, draw, thick, rounded corners, minimum width=3cm, minimum height=1cm, align=center, font=\small},
  container/.style={rectangle, draw, very thick, dashed, rounded corners, inner sep=0.4cm},
  arrow/.style={-Latex, thick},
  dasharrow/.style={-Latex, thick, dashed},
]
  \node[container, label={\textbf{deephold\_db -- Entwicklungsmaschine}}] (host) {
    \begin{tikzpicture}[node distance=0.6cm]
      \node[box, fill=blue!20] (py) {Python venv\\\scriptsize deephold\_db package};
      \node[box, fill=green!20, right=of py] (pg) {PostgreSQL 16\\\scriptsize in Docker};
      \node[box, fill=yellow!20, right=of pg] (bun) {Bun\\\scriptsize deephold-db-tui};
      \node[box, fill=cyan!20, below=of pg] (admin) {Adminer\\\scriptsize Port 8080};
      \node[box, fill=orange!20, right=of admin] (prefect) {Prefect Server\\\scriptsize Port 4200};
    \end{tikzpicture}
  };

  \draw[arrow] (py) -- node[above, font=\scriptsize] {psycopg / SQLAlchemy} (pg);
  \draw[arrow] (bun) -- node[above, font=\scriptsize] {node-postgres} (pg);
  \draw[arrow] (pg) -- (admin);
  \draw[arrow] (pg) -- (prefect);
\end{tikzpicture}
\caption{Lokale Infrastruktur: Vier Prozesse, eine Datenbank. Alles
  läuft auf einer Maschine. PostgreSQL ist der zentrale Hub; Python
  und Bun kommunizieren nicht direkt, sondern nur über die DB.}
\label{fig:infra-overview}
\end{figure}

\begin{definition}
\textbf{localhost / 127.0.0.1}
\emph{Localhost} ist der Rechnername, der immer auf den
\emph{aktuellen Rechner} selbst verweist. Die IP-Adresse~127.0.0.1
ist die \emph{Loopback}-Adresse. Wenn ein Prozess auf Localhost hört,
ist er nur für Prozesse auf derselben Maschine erreichbar -- nicht
für andere Computer im Netz. Dies ist die sicherste
\emph{Standardtopologie} für Entwicklung.
\end{definition}

\section{Docker und Docker Compose}

\begin{definition}
\textbf{docker-compose.yml}
Die Datei \code{docker-compose.yml} ist die
\emph{Single Source of Truth} für alle containerisierten Services.
Sie listet jeden Service mit Image, Umgebungsvariablen, Ports,
Volumes und Healthcheck auf.
\end{definition}

\begin{lstlisting}[language=yaml, caption={Vollständige \code{docker-compose.yml}},label=lst:dockercompose]
name: deephold-db

services:
  postgres:
    image: postgres:16-alpine
    container_name: deephold_pg
    restart: unless-stopped
    environment:
      POSTGRES_USER: ${POSTGRES_USER:-deephold}
      POSTGRES_PASSWORD: ${POSTGRES_PASSWORD:-deephold}
      POSTGRES_DB: ${POSTGRES_DB:-deephold}
    ports:
      - "${POSTGRES_PORT:-5432}:5432"
    volumes:
      - postgres_data:/var/lib/postgresql/data
    healthcheck:
      test: ["CMD-SHELL", "pg_isready -U ${POSTGRES_USER:-deephold} -d ${POSTGRES_DB:-deephold}"]
      interval: 5s
      timeout: 5s
      retries: 10

  adminer:
    image: adminer:4
    container_name: deephold_adminer
    restart: unless-stopped
    ports:
      - "${ADMINER_PORT:-8080}:8080"
    environment:
      ADMINER_DEFAULT_SERVER: postgres
    depends_on:
      postgres:
        condition: service_healthy

  prefect-server:
    image: prefecthq/prefect:2.16-python3.11
    container_name: deephold_prefect
    restart: unless-stopped
    command: prefect server start --host 0.0.0.0
    ports:
      - "${PREFECT_PORT:-4200}:4200"
    environment:
      PREFECT_HOME: /opt/prefect
      PREFECT_SERVER_API_URL: http://localhost:4200/api
    volumes:
      - prefect_data:/opt/prefect

volumes:
  postgres_data:
  prefect_data:
\end{lstlisting}

Listing~\ref{lst:dockercompose} zeigt die vollständige
\code{docker-compose.yml}. Bemerkenswert:

\begin{itemize}
  \item \textbf{Drei Services}: \code{postgres} (Datenbank),
        \code{adminer} (Web-UI für SQL), \code{prefect-server}
        (Workflow-Orchestrierung).
  \item \textbf{Zwei Volumes}: \code{postgres\_data} und
        \code{prefect\_data} für persistente Daten.
  \item \textbf{Healthcheck} auf Postgres -- der \code{adminer}-Service
        startet erst, wenn Postgres bereit ist (\code{condition:
        service\_healthy}).
  \item \textbf{Umgebungsvariablen} mit Defaults:
        \code{\$\{POSTGRES\_USER:-deephold\}} -- falls die Variable
        nicht gesetzt ist, wird der Default \code{deephold} verwendet.
  \item \textbf{Image-Tags}: \code{postgres:16-alpine} -- die
        \code{-alpine}-Variante nutzt eine minimale Linux-Distribution
        (Alpine) und ist entsprechend klein.
\end{itemize}

\section{Das Makefile}

Das Makefile ist die \emph{einzige CLI}, die man kennen muss. Es
abstrahiert über alle Details -- Docker, Python, Bun, Alembic,
Tests.

\begin{lstlisting}[language=make, caption={Vollständiges Makefile},label=lst:makefile]
.PHONY: help
help:  ## Zeige alle Targets
	@grep -E '^[a-zA-Z_-]+:.*?## .*$$' $(MAKEFILE_LIST) | \
	  awk 'BEGIN {FS = ":.*?## "}; {printf "\033[36m%-20s\033[0m %s\n", $$1, $$2}'

.PHONY: init
init:  ## Komplett-Setup: deps + Stack + Migrations
	poetry install
	$(COMPOSE) up -d
	@echo "Warte auf Postgres ..."
	@for i in $$(seq 1 30); do \
	  $(COMPOSE) exec -T postgres pg_isready -U $${POSTGRES_USER:-deephold} >/dev/null 2>&1 && break; \
	  sleep 1; \
	done
	$(MAKE) migrate
	@echo "Init fertig. UI: Adminer :8080, Prefect :4200"

.PHONY: up
up:  ## docker-compose Stack starten
	$(COMPOSE) up -d

.PHONY: down
down:  ## docker-compose Stack stoppen
	$(COMPOSE) down

.PHONY: migrate
migrate:  ## Alembic-Migrationen anwenden
	DATABASE_URL=$(DB_DSN) $(PYTHON) -m alembic upgrade head

.PHONY: revision
revision:  ## Neue Alembic-Revision (Usage: make revision m="...")
	DATABASE_URL=$(DB_DSN) $(PYTHON) -m alembic revision --autogenerate -m "$(m)"

.PHONY: test
test:  ## pytest + pandera schemas
	DATABASE_URL=$(DB_DSN) $(PYTHON) -m pytest

.PHONY: tui
tui:  ## OpenTUI-Explorer starten (Bun erforderlich)
	@if ! command -v bun >/dev/null 2>&1; then \
		echo "Bun nicht installiert. Install: curl -fsSL https://bun.sh/install | bash"; \
		exit 1; \
	fi
	cd tui && bun install --silent && bun run start

.PHONY: tui-test
tui-test:  ## TUI-Tests (Unit + Integration)
	@if ! command -v bun >/dev/null 2>&1; then \
		echo "Bun nicht installiert. Install: curl -fsSL https://bun.sh/install | bash"; \
		exit 1; \
	fi
	cd tui && bun test

.PHONY: format
format:  ## ruff format
	$(PYTHON) -m ruff format .

.PHONY: lint
lint:  ## ruff check
	$(PYTHON) -m ruff check .

.PHONY: typecheck
typecheck:  ## mypy
	$(PYTHON) -m mypy src/

.PHONY: check
check:  ## Alles: lint + typecheck + test
	$(PYTHON) -m ruff check .
	$(PYTHON) -m ruff format --check .
	$(PYTHON) -m mypy src/
	$(PYTHON) -m pytest
\end{lstlisting}

Listing~\ref{lst:makefile} zeigt die wichtigsten Targets. Die
Konvention ist: \code{make <target>} -- und das Makefile sucht
sich die richtigen Befehle zusammen.

\section{Python-Environment}

\begin{definition}
\textbf{venv (Virtual Environment)}
Ein \emph{virtuelles Python-Environment} ist ein isoliertes
Verzeichnis, in dem Python-Pakete installiert werden, ohne das
System-Python zu beeinflussen. In diesem Projekt wird statt
Poetry (das viele zusätzliche Komplexität mitbringt) direkt
\code{python -m venv .venv} und \code{pip install} verwendet -- so
funktioniert das Setup auf jeder Maschine ohne spezielle Tools.
\end{definition}

\begin{beispiel}
\textbf{Setup-Sequenz}
\begin{lstlisting}[language=shell]
# 1. Virtuelles Environment anlegen
python3 -m venv .venv

# 2. pip aktualisieren
.venv/bin/pip install --quiet --upgrade pip

# 3. Projekt installieren (inkl. dev-deps)
.venv/bin/pip install --quiet -e .

# 4. Tests laufen lassen
.venv/bin/pytest
\end{lstlisting}
\end{beispiel}

Die wichtigsten Python-Pakete sind in
Tabelle~\ref{tab:py-deps} zusammengefasst.

\begin{table}[ht]
\centering
\small
\begin{tabular}{lll}
  \toprule
  \textbf{Paket} & \textbf{Version} & \textbf{Zweck} \\
  \midrule
  \code{polars}    & $\geq$0.20 & Datenverarbeitung (schneller als Pandas) \\
  \code{pandas}    & $\geq$2.2  & Datenverarbeitung (Kompatibilität) \\
  \code{numpy}     & $\geq$1.26 & Numerik \\
  \code{sqlalchemy} & $\geq$2.0  & ORM \\
  \code{alembic}   & $\geq$1.13 & Schema-Migrationen \\
  \code{psycopg[binary]} & $\geq$3.1 & PostgreSQL-Treiber \\
  \code{httpx}      & $\geq$0.27 & HTTP-Client (Vendor-Adapter) \\
  \code{tenacity}   & $\geq$8.2  & Retries \\
  \code{structlog}  & $\geq$24.1 & Logging \\
  \code{pydantic}   & $\geq$2.6  & Datenmodelle \\
  \code{pydantic-settings} & $\geq$2.2 & Config aus .env \\
  \code{prefect}    & $\geq$2.16 & Workflow-Orchestrierung \\
  \code{tabulate}   & $\geq$0.9  & Formatierte Tabellen im Terminal \\
  \midrule
  \code{pytest}     & $\geq$8.0  & Test-Runner \\
  \code{ruff}       & $\geq$0.3  & Linter + Formatter \\
  \code{mypy}       & $\geq$1.8  & Type-Checker \\
  \bottomrule
\end{tabular}
\caption{Wichtigste Python-Abhängigkeiten}
\label{tab:py-deps}
\end{table}

\section{Bun und TypeScript}

Bun ist die bevorzugte Runtime für das TUI (Kapitel~\ref{ch:tui}).
Es vereint Paketmanager, Test-Runner und Bundler in einem einzigen,
sehr schnellen Binary.

\begin{lstlisting}[language=shell, caption={Bun-Setup}]
# Installation (einmalig pro Maschine)
curl -fsSL https://bun.sh/install | bash
source ~/.bashrc

# Im tui/-Verzeichnis
cd tui
bun install                  # Installiert Dependencies
bun run start                # Startet TUI
bun test                     # Lässt alle Tests laufen
bun run typecheck            # tsc --noEmit
\end{lstlisting}

\section{Pfade, Ports, Datenverzeichnisse}

\begin{table}[ht]
\centering
\small
\begin{tabular}{lll}
  \toprule
  \textbf{Was} & \textbf{Pfad / Port} & \textbf{Wer greift zu} \\
  \midrule
  Postgres-Server         & \code{localhost:5432}    & Python, Bun, Adminer, psql \\
  Adminer-Web-UI          & \code{localhost:8080}    & Browser \\
  Prefect-Server-UI       & \code{localhost:4200}    & Browser \\
  PostgreSQL-Daten-Volume & \code{postgres\_data}    & Docker (persistent) \\
  Prefect-Daten-Volume    & \code{prefect\_data}     & Docker (persistent) \\
  Python-Pakete           & \code{./.venv/}          & \code{./.venv/bin/python} \\
  Bun-Cache               & \code{$\sim$/.bun/}      & Bun \\
  TUI-Node-Modules        & \code{./tui/node\_modules/} & Bun \\
  Jupyter-Notebook-Output & \code{./notebooks/}      & Browser (HTML) \\
  \bottomrule
\end{tabular}
\caption{Pfade und Ports aller Komponenten}
\label{tab:paths-ports}
\end{table}

\section{Was beim ersten Start passiert}

\begin{example}[Boot-Sequenz]
\begin{enumerate}
  \item \code{make up} startet die Docker-Container.
  \item Postgres initialisiert sich (5--10 Sekunden) und ist bereit
        für Verbindungen (\code{pg\_isready} gibt Exit-Code 0).
  \item Adminer wartet auf \code{postgres: service\_healthy},
        startet dann.
  \item Prefect-Server startet (langsamer, 10--20 Sekunden).
  \item \code{make migrate} führt alle Alembic-Migrationen aus und
        erzeugt die 14~Tabellen.
  \item \code{make ingest-<asset\_class>} oder \code{pytest} oder
        \code{make tui} können jetzt starten.
\end{enumerate}
\end{example}

\section{Was als Nächstes?}

Mit der Infrastruktur im Griff wenden wir uns dem nächsten
zentralen Thema zu: dem \emph{Datenmodell}. Kapitel~\ref{ch:datenmodell}
erklärt die 14~Tabellen, ihre Beziehungen und ihre Indizes.

% =====================================================================
% Kapitel 4: Datenmodell
% =====================================================================

\chapter{Datenmodell}

\label{ch:datenmodell}

Das Datenmodell ist das \emph{Fundament} des gesamten Projekts.
Es definiert, welche Daten gespeichert werden, wie sie
zusammenhängen und wie effizient wir sie abfragen können.

\section{Überblick: Die 14~Tabellen}

Tabelle~\ref{tab:tables} gibt einen schnellen Überblick über alle
14~Tabellen. Die vollständigen DDL-Definitionen finden sich in
Anhang~C.

\begin{table}[ht]
\centering
\small
\begin{tabular}{llll}
  \toprule
  \textbf{Tabelle} & \textbf{Zweck} & \textbf{PK} & \textbf{Wichtige FKs} \\
  \midrule
  \code{venues}              & Börsen/Handelsplätze       & \code{venue\_id}     & -- \\
  \code{vendors}            & Datenlieferanten           & \code{vendor\_id}    & -- \\
  \code{instruments}        & Asset-Master               & \code{instrument\_id} & \code{primary\_venue} \\
  \code{instrument\_identifiers} & Cross-Vendor-Mapping   & (\code{scheme}, \code{value}) & \code{instrument\_id} \\
  \code{trading\_calendars} & Handelstage je Venue       & (\code{venue\_id}, \code{date}) & \code{venue\_id} \\
  \code{prices\_raw}        & Roh-Preise (vendor-original) & (\code{vendor\_id}, \code{vendor\_symbol}, \code{date}) & \code{vendor\_id} \\
  \code{prices\_daily}      & OHLCV (harmonisiert)      & (\code{instrument\_id}, \code{date}) & \code{instrument\_id}, \code{vendor\_id} \\
  \code{corporate\_actions} & Splits, Dividenden         & (\code{instrument\_id}, \code{ex\_date}, \code{action\_type}) & \code{instrument\_id} \\
  \code{fx\_rates\_daily}   & Cross-Currency-FX          & (\code{ccy\_from}, \code{ccy\_to}, \code{date}) & \code{vendor\_id} \\
  \code{bond\_yields}        & Bond-Yield-Curve           & (\code{instrument\_id}, \code{date}) & \code{instrument\_id} \\
  \code{macro\_series}      & Makro-Stammdaten           & \code{series\_id}    & -- \\
  \code{macro\_observations}& Makro-Beobachtungen        & (\code{series\_id}, \code{date}) & \code{series\_id} \\
  \code{dq\_runs}            & DQ-Lauf-Metadaten          & \code{run\_id}       & -- \\
  \code{dq\_findings}        & DQ-Befunde                 & (\code{run\_id}, \code{check\_name}) & \code{run\_id} \\
  \bottomrule
\end{tabular}
\caption{Alle 14~Tabellen des \code{deephold\_db}-Schemas}
\label{tab:tables}
\end{table}

\section{ER-Diagramm}

Abbildung~\ref{fig:er} zeigt das vollständige ER-Diagramm.

\begin{figure}[ht]
\centering
\resizebox{\textwidth}{!}{%
\begin{tikzpicture}[
  table/.style={
    rectangle, draw, thick, rounded corners=2pt,
    minimum width=4cm, align=left, inner sep=4pt, font=\scriptsize,
  },
  title/.style={font=\bfseries, fill=blue!20, anchor=west, inner sep=2pt},
  col/.style={font=\ttfamily},
  pk/.style={col, fill=yellow!20},
  fk/.style={col, fill=green!20},
  arrow/.style={-Latex, thick},
]
  % instruments
  \node[table] (instruments) {
    \begin{tabular}{l}
      \node[title]{instruments}; \\
      \node[pk]{PK instrument\_id}; \\
      \node[col]{asset\_class (TEXT, NN)}; \\
      \node[col]{name (TEXT, NN)}; \\
      \node[col]{currency (CHAR(3))}; \\
      \node[fk]{FK primary\_venue $\to$ venues}; \\
      \node[col]{is\_active (BOOL)}; \\
    \end{tabular}
  };

  % venues
  \node[table, left=2.5cm of instruments] (venues) {
    \begin{tabular}{l}
      \node[title]{venues}; \\
      \node[pk]{PK venue\_id}; \\
      \node[col]{code (TEXT, UQ)}; \\
      \node[col]{name (TEXT)}; \\
      \node[col]{timezone (TEXT)}; \\
      \node[col]{country (CHAR(2))}; \\
    \end{tabular}
  };

  % vendors
  \node[table, above=1.2cm of instruments] (vendors) {
    \begin{tabular}{l}
      \node[title]{vendors}; \\
      \node[pk]{PK vendor\_id}; \\
      \node[col]{code (TEXT, UQ)}; \\
      \node[col]{license (TEXT)}; \\
      \node[col]{homepage (TEXT)}; \\
    \end{tabular}
  };

  % instrument_identifiers
  \node[table, right=2.5cm of instruments] (identifiers) {
    \begin{tabular}{l}
      \node[title]{instrument\_identifiers}; \\
      \node[pk]{PK (scheme, value)}; \\
      \node[fk]{FK instrument\_id $\to$ instruments}; \\
      \node[col]{scheme (TEXT, NN)}; \\
      \node[col]{value (TEXT, NN)}; \\
    \end{tabular}
  };

  % prices_daily
  \node[table, below=1.2cm of instruments] (prices) {
    \begin{tabular}{l}
      \node[title]{prices\_daily}; \\
      \node[pk]{PK (instrument\_id, date)}; \\
      \node[fk]{FK instrument\_id $\to$ instruments}; \\
      \node[fk]{FK vendor\_id $\to$ vendors}; \\
      \node[col]{open, high, low, close, adj\_close}; \\
      \node[col]{volume (BIGINT)}; \\
    \end{tabular}
  };

  % prices_raw
  \node[table, left=2.5cm of prices] (pricesraw) {
    \begin{tabular}{l}
      \node[title]{prices\_raw}; \\
      \node[pk]{PK (vendor\_id, vendor\_symbol, date)}; \\
      \node[fk]{FK vendor\_id $\to$ vendors}; \\
      \node[col]{payload (JSONB, NN)}; \\
    \end{tabular}
  };

  % macro_series
  \node[table, right=2.5cm of prices] (macroseries) {
    \begin{tabular}{l}
      \node[title]{macro\_series}; \\
      \node[pk]{PK series\_id (TEXT)}; \\
      \node[col]{name, source, frequency, unit}; \\
    \end{tabular}
  };

  % macro_observations
  \node[table, below=1.2cm of prices] (macroobs) {
    \begin{tabular}{l}
      \node[title]{macro\_observations}; \\
      \node[pk]{PK (series\_id, date)}; \\
      \node[fk]{FK series\_id $\to$ macro\_series}; \\
      \node[col]{value (NUMERIC(18,6))}; \\
    \end{tabular}
  };

  % FX
  \node[table, right=2.5cm of macroobs] (fx) {
    \begin{tabular}{l}
      \node[title]{fx\_rates\_daily}; \\
      \node[pk]{PK (ccy\_from, ccy\_to, date)}; \\
      \node[col]{rate (NUMERIC(18,8))}; \\
      \node[fk]{FK vendor\_id $\to$ vendors}; \\
    \end{tabular}
  };

  % corporate_actions
  \node[table, left=2.5cm of macroobs] (corpactions) {
    \begin{tabular}{l}
      \node[title]{corporate\_actions}; \\
      \node[pk]{PK (instrument\_id, ex\_date, action\_type)}; \\
      \node[fk]{FK instrument\_id $\to$ instruments}; \\
      \node[col]{value (NUMERIC(18,8))}; \\
    \end{tabular}
  };

  % bond_yields
  \node[table, right=2.5cm of macroobs] (bondyields) {
    \begin{tabular}{l}
      \node[title]{bond\_yields}; \\
      \node[pk]{PK (instrument\_id, date)}; \\
      \node[fk]{FK instrument\_id $\to$ instruments}; \\
      \node[col]{yield, price, duration, rating}; \\
    \end{tabular}
  };

  % dq
  \node[table, below=1.5cm of macroobs] (dq) {
    \begin{tabular}{l}
      \node[title]{dq\_runs / dq\_findings}; \\
      \node[col]{Lauf-Metadaten + Befunde}; \\
    \end{tabular}
  };

  % Verbindungen
  \draw[arrow] (venues) -- (instruments);
  \draw[arrow] (vendors) -- (prices);
  \draw[arrow] (instruments) -- (identifiers);
  \draw[arrow] (instruments) -- (prices);
  \draw[arrow] (instruments) -- (corpactions);
  \draw[arrow] (instruments) -- (bondyields);
  \draw[arrow] (macroseries) -- (macroobs);
\end{tikzpicture}%
}
\caption{Entity-Relationship-Diagramm des \code{deephold\_db}-Schemas.
  \code{instruments} ist der zentrale Asset-Master; \code{vendors} ist
  die Quelle aller Preise. Makro- und Equity-Daten sind
  \emph{parallel} organisiert, nicht in einer Hierarchie.}
\label{fig:er}
\end{figure}

\section{Stammdaten-Tabellen}

\subsection{\code{venues}}

\emph{Börsen und Handelsplätze}. Beispiele: NASDAQ, NYSE, XETRA,
LSE, TYO. Jede Zeile beschreibt einen Handelsplatz eindeutig über
den Code (\code{NASDAQ}, \code{NYSE}, \code{XETRA} \ldots). Die
\code{timezone} ermöglicht später korrekte Datums-Berechnungen über
Zeitzonen hinweg.

\subsection{\code{vendors}}

\emph{Datenlieferanten} -- die Quellen, aus denen wir Preise
beziehen. Beispiele: \code{fred}, \code{ecb}, \code{yahoo}, \code{demo}.
Die \code{license}-Spalte speichert die Nutzungsbedingungen (z.\,B.\
\foreign{public domain}, \foreign{CC BY 4.0}, \foreign{Yahoo ToS --
privater Gebrauch}).

\subsection{\code{instruments}}

\emph{Der Asset-Master} -- eine zentrale Tabelle, in der jedes
gehandelte Wertpapier genau einmal vorkommt: eine Aktie, ein ETF,
ein Index, eine Anleihe. \code{asset\_class} ist ein
\emph{Diskriminator} (\code{equity}, \code{index}, \code{bond\_gov},
\ldots), der bestimmt, in welchen weiteren Tabellen das Instrument
referenziert wird.

\subsection{\code{instrument\_identifiers}}

\emph{Das Cross-Vendor-Mapping}. Ein Instrument kann unter
verschiedenen IDs bei verschiedenen Vendoren bekannt sein:

\begin{itemize}
  \item Yahoo Finance: \code{AAPL}
  \item ISIN: \code{US0378331005}
  \item Bloomberg: \code{AAPL:US}
  \item FRED: \code{DEXUSEU} (für EUR/USD-FX)
\end{itemize}

\code{instrument\_identifiers} bildet jede dieser IDs auf die
\emph{interne} \code{instrument\_id} ab. Der zusammengesetzte
Primärschlüssel (\code{scheme}, \code{value}) verhindert Duplikate.

\section{Preis-Tabellen}

\subsection{\code{prices\_raw} -- Das unveränderte Original}

\code{prices\_raw} speichert die \emph{exakten Roh-Payloads} der
Vendoren, ohne Interpretation. Jede Zeile enthält:

\begin{itemize}
  \item \code{vendor\_id} -- welcher Lieferant
  \item \code{vendor\_symbol} -- dessen nativers Ticker-Symbol
  \item \code{date} -- Handelstag
  \item \code{payload} -- die vollständige rohe Antwort als JSONB
  \item \code{ingested\_at} -- Zeitstempel des Pulls
\end{itemize}

\begin{merke}
\code{prices\_raw} ist der \emph{Append-Only Audit-Trail}. Wir
schreiben hinein, ändern aber nie. Selbst wenn die harmonisierte
Tabelle \code{prices\_daily} korrigiert werden muss, ist die
Original-Antwort noch da.
\end{merke}

\subsection{\code{prices\_daily} -- Die harmonisierte Zeitreihe}

\code{prices\_daily} ist die \emph{Arbeitstabelle} für alle
Preis-basierten Analysen. Sie ist auf \code{instrument\_id}
\emph{normalisiert} (nicht auf das Vendor-Symbol) -- dadurch
können wir Daten verschiedener Vendoren für dasselbe Instrument
zusammenführen und vergleichen.

\begin{lstlisting}[language=sql, caption={Auszug aus \code{prices\_daily}}]
\d prices_daily
\end{lstlisting}

\begin{lstlisting}[language=sql]
                                       Table "public.prices_daily"
       Column        |           Type           | Nullable |    Default
--------------------+--------------------------+----------+---------------
 instrument_id       | bigint                   | not null |
 date                | date                     | not null |
 open                | numeric(18,8)            |          |
 high                | numeric(18,8)            |          |
 low                 | numeric(18,8)            |          |
 close               | numeric(18,8)            |          |
 adjusted_close      | numeric(18,8)            |          |
 volume              | bigint                   |          |
 vendor_id           | integer                  |          |
Indexes:
    "prices_daily_pkey" PRIMARY KEY, btree (instrument_id, date)
    "ix_prices_daily_date" btree (date)
    "ix_prices_daily_instrument_id_date" btree (instrument_id, date)
Check constraints:
    "prices_daily_check" CHECK (instrument_id IS NOT NULL)
Foreign-key constraints:
    "prices_daily_instrument_id_fkey" FOREIGN KEY (instrument_id) REFERENCES instruments(instrument_id)
    "prices_daily_vendor_id_fkey" FOREIGN KEY (vendor_id) REFERENCES vendors(vendor_id)
\end{lstlisting}

Beachten Sie die Indizes: \code{ix\_prices\_daily\_date} beschleunigt
\emph{zeitbasierte} Queries (\code{WHERE date BETWEEN ...}), der
zusammengesetzte Index \code{ix\_prices\_daily\_instrument\_id\_date}
beschleunigt Queries nach \emph{einem Instrument über die Zeit}.

\section{Makro-Tabellen}

\subsection{\code{macro\_series} und \code{macro\_observations}}

Makro-Daten (Zinssätze, Inflation, Arbeitsmarkt) werden in zwei
\emph{getrennte} Tabellen aufgeteilt:

\begin{itemize}
  \item \code{macro\_series} -- Stammdaten: Name, Quelle, Frequenz,
        Einheit.
  \item \code{macro\_observations} -- die eigentlichen Werte, eine
        Zeile pro (Serie, Datum).
\end{itemize}

Diese Aufteilung ist ein klassisches
\emph{Star-Schema}: Die \code{series}-Tabelle ist die
\emph{Dimension}, \code{observations} ist die \emph{Fact}-Tabelle.
Sie ermöglicht, eine Serie nachträglich um Metadaten zu ergänzen
(z.\,B.\ um die Datenquelle zu wechseln), ohne alle Beobachtungen
anzufassen.

\section{FX- und Bond-Tabellen}

\subsection{\code{fx\_rates\_daily}}

Wechselkurse -- eine Zeile pro (\code{ccy\_from}, \code{ccy\_to},
\code{date}). Die \code{ccy\_from}/\code{ccy\_to}-Codierung
statt einer \code{rate}-Tabelle mit zwei Zeilen pro Paar spart
Speicher und macht Queries schneller.

\subsection{\code{bond\_yields}}

Bond-spezifische Daten -- die allgemeine \code{prices\_daily}-Tabelle
enthält keine Yield-Curve-Daten. Hier werden zusätzlich
\code{yield}, \code{price}, \code{duration} und \code{rating}
gespeichert.

\section{DQ-Tabellen}

\subsection{\code{dq\_runs} und \code{dq\_findings}}

\emph{Data-Quality}-Tabellen. \code{dq\_runs} protokolliert jeden
Lauf eines DQ-Checks (z.\,B.\ \foreign{OHLC-Konsistenz} oder
\foreign{no negative prices}), \code{dq\_findings} die einzelnen
Befunde. Diese Tabellen sind derzeit nur \emph{Skelette} -- die
konkreten Checks werden in einem späteren Kapitel erläutert.

\section{SQLAlchemy-Definitionen}

Das Schema wird in Python deklariert. Listing~\ref{lst:models}
zeigt einen Auszug.

\begin{lstlisting}[language=python, caption={Auszug aus \code{src/deephold\_db/db/models.py}}]
class Instrument(Base):
    __tablename__ = "instruments"
    instrument_id: Mapped[int] = mapped_column(BigInteger, primary_key=True, autoincrement=True)
    asset_class: Mapped[str] = mapped_column(Text, nullable=False)
    name: Mapped[str] = mapped_column(Text, nullable=False)
    currency: Mapped[str | None] = mapped_column(CHAR(3))
    primary_venue: Mapped[int | None] = mapped_column(Integer, ForeignKey("venues.venue_id"))
    is_active: Mapped[bool] = mapped_column(Boolean, server_default=func.true())
    created_at: Mapped[datetime] = mapped_column(DateTime(timezone=True), server_default=func.now())

class PricesDaily(Base):
    __tablename__ = "prices_daily"
    instrument_id: Mapped[int] = mapped_column(
        BigInteger, ForeignKey("instruments.instrument_id"), primary_key=True
    )
    date: Mapped[date] = mapped_column(Date, primary_key=True)
    open: Mapped[Decimal | None] = mapped_column(Numeric(18, 8))
    high: Mapped[Decimal | None] = mapped_column(Numeric(18, 8))
    low: Mapped[Decimal | None] = mapped_column(Numeric(18, 8))
    close: Mapped[Decimal | None] = mapped_column(Numeric(18, 8))
    adjusted_close: Mapped[Decimal | None] = mapped_column(Numeric(18, 8))
    volume: Mapped[int | None] = mapped_column(BigInteger)
    vendor_id: Mapped[int | None] = mapped_column(Integer, ForeignKey("vendors.vendor_id"))
    __table_args__ = (
        Index("ix_prices_daily_date", "date"),
        Index("ix_prices_daily_instrument_id_date", "instrument_id", "date"),
    )
\end{lstlisting}

Listing~\ref{lst:models} zeigt die zentrale
\emph{Declarative}-Syntax von SQLAlchemy~2.x. Die
\code{Mapped[]}-Annotation gibt Typ + Optionalität an, der
\code{mapped\_column()}-Aufruf definiert die
Datenbank-Repräsentation.

\section{Was als Nächstes?}

Mit dem Schema im Hinterkopf wenden wir uns den \emph{Vendor-Adaptern}
zu -- den Komponenten, die Daten aus FRED, ECB und Yahoo Finance
in dieses Schema einfüllen. Kapitel~\ref{ch:vendor}.

% =====================================================================
% Kapitel 5: Vendor-Adapter
% =====================================================================

\chapter{Vendor-Adapter}

\label{ch:vendor}

\emph{Vendor-Adapter} sind die Komponenten, die mit den externen
Datenquellen sprechen. Jeder Adapter hat die gleiche Schnittstelle
-- \code{fetch(symbol, start, end)} und \code{healthcheck()} -- aber
unterschiedliche Implementierungen je nach API-Protokoll der
Datenquelle.

\section{Das Vendor-Interface}

\begin{definition}
\textbf{Vendor-Adapter}
Ein \emph{Vendor-Adapter} ist eine Python-Klasse, die eine bestimmte
externe Datenquelle (z.\,B.\ FRED) abstrahiert. Sie bietet eine
\emph{einheitliche} Schnittstelle, damit der Rest des Systems nicht
wissen muss, ob die Daten von einer REST-API, einer SDMX-Quelle
oder einer CSV-Datei kommen.
\end{definition}

Listing~\ref{lst:vendor-base} zeigt die
\emph{Abstrakte-Basisklasse}, von der alle konkreten Adapter erben.

\begin{lstlisting}[language=python, caption={Basisklasse aller Vendor-Adapter},label=lst:vendor-base]
from abc import ABC, abstractmethod
from datetime import date
import polars as pl

class Vendor(ABC):
    """Base class for all vendor adapters."""
    code: str              # e.g. "fred", "ecb", "yahoo"
    base_url: str          # Vendor's API base URL

    @abstractmethod
    def fetch(self, symbol: str, start: date, end: date) -> pl.DataFrame:
        """Fetch raw data for `symbol` in [start, end] inclusive."""
        raise NotImplementedError

    @abstractmethod
    def healthcheck(self) -> bool:
        """Return True if the vendor is reachable."""
        raise NotImplementedError
\end{lstlisting}

Die Wahl von \code{polars.DataFrame} als Rückgabetyp ist
bewusst: Polars ist schnell, hat Lazy-Evaluation und ist vollständig
typisiert. Die Adapter sind dünn -- sie übersetzen JSON oder XML in
ein DataFrame -- und die schwere Arbeit (Statistiken, Resampling,
Joins) passiert downstream in dedizierten Analytics-Modulen
(siehe Kapitel~\ref{ch:analytics}).

\section{FRED -- Federal Reserve Economic Data}

\begin{definition}
\textbf{FRED}
\emph{FRED} (Federal Reserve Economic Data) ist die
Wirtschaftsdatenbank der US-Notenbank Federal Reserve. Sie enthält
über 800.000~Zeitreihen zu Zinssätzen, Inflation, Beschäftigung,
BIP, Geldmenge und vielem mehr. Die Daten sind \foreign{public domain}
und über eine kostenlose REST-API zugänglich.
\end{definition}

\begin{definition}
\textbf{API-Key}
Ein \emph{API-Key} ist ein geheimer Token, der den
\emph{API-Aufrufer} identifiziert. FRED verlangt einen Key, den man
kostenlos unter \url{https://fred.stlouisfed.org/docs/api/api_key.html}
beantragen kann. Der Key wird in der Datei \code{.env} (nie im
Quellcode!) gespeichert.
\end{definition}

FRED liefert Zeitreihen über die Endpoints
\code{/fred/series/observations} (für eine einzelne Zeitreihe) und
\code{/fred/series} (für Metadaten). Listing~\ref{lst:fred-fetch}
zeigt, wie ein typischer Fetch aussieht.

\begin{lstlisting}[language=python, caption={FRED-Vendor: Healthcheck und Fetch},label=lst:fred-fetch]
import httpx
from deephold_db.utils.rate_limit import limit
from deephold_db.utils.retry import http_retry
from deephold_db.utils.config import get_settings

class FredVendor(Vendor):
    code = "fred"
    base_url = "https://api.stlouisfed.org/fred"

    def __init__(self, api_key: str | None = None) -> None:
        self.api_key = api_key or get_settings().fred_api_key
        if not self.api_key:
            raise ValueError("FRED_API_KEY is required (set in .env)")

    @http_retry()
    def _get(self, path: str, params: dict) -> dict:
        with limit(self.code, 120, 30):  # 120 req/min, burst 30
            q = {**params, "api_key": self.api_key, "file_type": "json"}
            r = httpx.get(f"{self.base_url}{path}", params=q, timeout=30.0)
            r.raise_for_status()
            return r.json()

    def healthcheck(self) -> bool:
        try:
            data = self._get("/releases", {"limit": 1})
            return "releases" in data
        except Exception:
            return False

    def fetch(self, symbol: str, start, end) -> pl.DataFrame:
        data = self._get(
            "/series/observations",
            {"series_id": symbol, "observation_start": start.isoformat(),
             "observation_end": end.isoformat()},
        )
        rows = [{"date": obs["date"], "value": _parse_value(obs["value"])}
                for obs in data.get("observations", [])]
        return pl.DataFrame(rows, schema={"date": pl.Date, "value": pl.Float64})
\end{lstlisting}

Beachten Sie:

\begin{itemize}
  \item \code{@http\_retry()} -- dekoriert die Methode mit automatischen
        Retries bei Netzwerkfehlern (3 Versuche, 5s/30s/120s Backoff).
  \item \code{with limit(...)} -- Token-Bucket, das 120~req/min
        erlaubt (FREDs Limit) mit Burst~30.
  \item \code{\_parse\_value} -- FRED verwendet \code{"."} für fehlende
        Werte; das wird zu \code{NaN} konvertiert.
\end{itemize}

\section{ECB -- European Central Bank SDMX}

\begin{definition}
\textbf{SDMX (Statistical Data and Metadata Exchange)}
\emph{SDMX} ist ein ISO-Standard für den Austausch statistischer Daten.
Große Institutionen wie EZB, OECD, IWF und Eurostat veröffentlichen
ihre Daten über SDMX-Endpunkte. Das Format hat zwei Häufigkeiten:

\begin{itemize}
  \item \textbf{SDMX-ML} (XML-basiert) -- der offizielle Standard
  \item \textbf{SDMX-JSON} (JSON-basiert) -- moderner, einfacher
        zu parsen
\end{itemize}

Dieses Projekt verwendet die \emph{CSV}-Variante (über
\code{format=csvdata}), weil sie sich am einfachsten in Polars laden
lässt.
\end{definition}

\begin{beispiel}
\textbf{ECB SDMX URL-Struktur}
\begin{lstlisting}[language=shell]
# EXR = Exchange Rates
# D.USD.EUR.SP00.A = Daily, USD, EUR, Spot, Average
https://data-api.ecb.europa.eu/service/data/EXR/D.USD.EUR.SP00.A?startPeriod=2024-01-01&endPeriod=2024-12-31&format=csvdata
\end{lstlisting}
\end{beispiel}

Die ECB-API ist im Gegensatz zu FRED \emph{anonym} -- kein API-Key
nötig. Listing~\ref{lst:ecb-fetch} zeigt den ECB-Vendor.

\begin{lstlisting}[language=python, caption={ECB-Vendor: Parse SDMX CSV},label=lst:ecb-fetch]
class ECBVendor(Vendor):
    code = "ecb"

    def __init__(self, base_url: str | None = None) -> None:
        self.base_url = (base_url or get_settings().ecb_sdmx_base).rstrip("/")

    @http_retry()
    def _get_csv(self, path: str, params: dict) -> str:
        with limit(self.code, 60, 10):
            q = {**params, "format": "csvdata"}
            r = httpx.get(f"{self.base_url}{path}", params=q, timeout=30.0)
            r.raise_for_status()
            return r.text

    def fetch(self, symbol, start, end) -> pl.DataFrame:
        flow, key = self._resolve_symbol(symbol)
        csv_text = self._get_csv(
            f"/data/{flow}/{key}",
            {"startPeriod": start.isoformat(), "endPeriod": end.isoformat()},
        )
        return _parse_ecb_csv(csv_text)
\end{lstlisting}

Die SDMX-Daten haben \emph{zwei verschiedene Datumsformate}:
Tagesdaten verwenden \code{YYYY-MM-DD} (z.\,B.\ \code{2024-06-03}),
Monatsdaten verwenden \code{YYYY-MM} (z.\,B.\ \code{2024-01}).
Der Parser muss beide Formen handhaben.

\begin{lstlisting}[language=python, caption={ECB-CSV-Parser: Beide Datumsformate}]
def _parse_ecb_csv(csv_text: str) -> pl.DataFrame:
    # Disables auto-type-inference (which would mis-parse "2024-01" as broken date)
    df = pl.read_csv(
        csv_text.encode("utf-8"),
        columns=["TIME_PERIOD", "OBS_VALUE"],
        dtypes={"TIME_PERIOD": pl.Utf8, "OBS_VALUE": pl.Float64},
    )
    if df.is_empty():
        return pl.DataFrame(schema={"date": pl.Date, "value": pl.Float64})

    # coalesce: first try full date, fall back to year-month
    df = df.with_columns(
        pl.coalesce(
            pl.col("TIME_PERIOD").str.strptime(pl.Date, format="%Y-%m-%d", strict=False),
            pl.col("TIME_PERIOD").str.strptime(pl.Date, format="%Y-%m", strict=False),
        ).alias("date")
    )
    return df.select(
        pl.col("date"),
        pl.col("OBS_VALUE").cast(pl.Float64, strict=False).alias("value"),
    )
\end{lstlisting}

\section{Yahoo Finance (via yfinance)}

\begin{definition}
\textbf{yfinance}
\emph{yfinance} ist eine Python-Bibliothek, die
\emph{inoffizielle} Yahoo-Finance-API verwendet, um historische
Aktien-, ETF- und Index-Daten abzurufen. Es ist \emph{nicht} offiziell
von Yahoo unterstützt -- die Bibliothek scraped und parst die
Yahoo-Website. Daher kann sie bei Yahoo-Änderungen kaputt gehen,
und die Yahoo \emph{terms of service} (ToS) erlauben
\emph{keine Weiterverteilung} der Daten. Dieses Projekt
verwendet yfinance daher nur für \emph{private} Forschungszwecke.
\end{definition}

\begin{beispiel}
\textbf{yfinance API -- sehr einfach}
\begin{lstlisting}[language=python]
import yfinance as yf
t = yf.Ticker("AAPL")
df = t.history(start="2024-01-01", end="2024-12-31", auto_adjust=False)
# df hat Spalten: Open, High, Low, Close, Volume, Dividends, Stock Splits
# Index: DatetimeIndex mit TZ (z.B. America/New_York)
\end{lstlisting}
\end{beispiel}

\begin{lstlisting}[language=python, caption={Yahoo-Vendor: Single und Bulk Fetch},label=lst:yahoo-fetch]
import yfinance as yf
import pandas as pd

class YahooVendor(Vendor):
    code = "yahoo"

    def fetch(self, symbol, start, end) -> pl.DataFrame:
        with limit(self.code, 30, 5):  # 30 req/min, burst 5
            t = yf.Ticker(symbol)
            pdf = t.history(
                start=start.isoformat(),
                end=(end + timedelta(days=1)).isoformat(),  # yf end is exclusive
                auto_adjust=False,    # we need both close and adj_close
                actions=False,         # splits/dividends go to corporate_actions
            )
        if pdf.empty:
            return pl.DataFrame(schema=_YAHOO_SCHEMA)
        return _yahoo_to_polars(pdf)

    def fetch_many(self, symbols: list[str], start, end) -> dict[str, pl.DataFrame]:
        """Bulk fetch in a single yfinance call -- much faster."""
        with limit(self.code, 30, 5):
            pdf = yf.download(
                tickers=" ".join(symbols),
                start=start.isoformat(),
                end=(end + timedelta(days=1)).isoformat(),
                auto_adjust=False,
                actions=False,
                group_by="ticker",
                threads=True,
                progress=False,
            )
        # ... handle MultiIndex columns for multiple tickers
\end{lstlisting}

Listing~\ref{lst:yahoo-fetch} zeigt den Yahoo-Vendor. Zwei
Dinge sind erwähnenswert:

\begin{enumerate}
  \item \textbf{Bulk-Methode}: \code{fetch\_many} verwendet
        \code{yf.download} mit \code{threads=True} -- dies holt
        \emph{alle} Symbole in einem einzigen API-Call. Für 36~Symbole
        benötigt es nur ~12~Sekunden, statt 36~einzelne Calls.
  \item \textbf{Pandas-zu-Polars-Bridge}: yfinance liefert ein
        \code{pandas.DataFrame} mit TZ-aware-DatetimeIndex und
        nullable Integers (\code{Int64} statt \code{int64}). Wir
        konvertieren explizit zu pl-kompatiblen Typen (\code{int64})
        und parsen die TZ-Datums in Polars.
\end{enumerate}

\section{Rate-Limiting und Retry}

\begin{definition}
\textbf{Exponential Backoff}
\emph{Exponential Backoff} ist eine Retry-Strategie, bei der die
Wartezeit zwischen Versuchen \emph{exponentiell wächst}: 5s, dann
30s, dann 120s. Dies gibt einem überlasteten Server Zeit, sich zu
erholen, ohne dass der Client ihn weiter bombardiert.
\end{definition}

Listing~\ref{lst:retry} zeigt die zentralen Werkzeuge für robustes
HTTP.

\begin{lstlisting}[language=python, caption={Rate-Limiting und Retry-Logik},label=lst:retry]
import httpx
import time
from tenacity import retry, stop_after_attempt, wait_chain, wait_fixed

# Token-Bucket für Rate-Limiting
class TokenBucket:
    def __init__(self, rate_per_min: int, burst: int):
        self.rate = rate_per_min / 60.0  # tokens per second
        self.burst = burst
        self.tokens = float(burst)
        self.last = time.monotonic()

    def take(self, n: int = 1) -> None:
        # Wenn nicht genug Tokens: blockiere bis genug da sind
        now = time.monotonic()
        elapsed = now - self.last
        self.last = now
        self.tokens = min(self.burst, self.tokens + elapsed * self.rate)
        if self.tokens < n:
            time.sleep((n - self.tokens) / self.rate)
            self.tokens = 0
        else:
            self.tokens -= n

# Retry-Decorator mit 5s/30s/120s Backoff
def http_retry():
    return retry(
        reraise=True,
        stop=stop_after_attempt(3),
        wait=wait_chain(wait_fixed(5), wait_fixed(30), wait_fixed(120)),
        retry=retry_if_exception_type((
            httpx.TimeoutException, httpx.NetworkError, httpx.RemoteProtocolError,
        )),
    )
\end{lstlisting}

Jeder Vendor-Adapter verwendet \code{limit(code, rate, burst)} und
\code{@http\_retry()}, um die spezifischen Anforderungen der Quelle
zu erfüllen. Tabelle~\ref{tab:rate-limits} zeigt die Limits pro
Vendor.

\begin{table}[ht]
\centering
\small
\begin{tabular}{lrrl}
  \toprule
  \textbf{Vendor} & \textbf{Rate} & \textbf{Burst} & \textbf{Retry-Backoff} \\
  \midrule
  FRED   & 120 req/min & 30 & 5s/30s/120s \\
  ECB    & 60 req/min  & 10 & 5s/30s/120s \\
  Yahoo  & 30 req/min  & 5  & 5s/30s/120s \\
  \bottomrule
\end{tabular}
\caption{Rate-Limits und Burst pro Vendor. Diese sind in
  \code{src/deephold\_db/vendors/<name>.py} konfiguriert.}
\label{tab:rate-limits}
\end{table}

\section{Was als Nächstes?}

Die Vendor-Adapter sind die Datenquellen; die ETL-Pipeline
(Kapitel~\ref{ch:etl}) zeigt, wie diese Daten in unsere
Datenbank gelangen und für Analysen verfügbar gemacht werden.

% =====================================================================
% Kapitel 6: ETL-Pipeline
% =====================================================================

\chapter{ETL-Pipeline}

\label{ch:etl}

Die \emph{ETL-Pipeline} -- Extract, Transform, Load -- ist das
\emph{Rückgrat} des Projekts. Sie ist verantwortlich dafür, dass
die Daten der Vendor-Adapter tatsächlich in die Datenbank gelangen
und für Analysen verfügbar werden.

\section{Überblick}

\begin{definition}
\textbf{ETL}
\emph{ETL} steht für:
  \begin{itemize}
    \item \textbf{Extract} -- Daten aus den Quellen (FRED, ECB, Yahoo)
          abrufen.
    \item \textbf{Transform} -- Daten in das richtige Format und Schema
          bringen.
    \item \textbf{Load} -- Daten in die Datenbank schreiben.
  \end{itemize}
\end{definition}

Abbildung~\ref{fig:etl-flow} zeigt den ETL-Flow im Detail.

\begin{figure}[ht]
\centering
\begin{tikzpicture}[
  node distance=0.8cm and 1.2cm,
  box/.style={rectangle, draw, thick, rounded corners, minimum width=3cm, minimum height=0.8cm, align=center, font=\small},
  arrow/.style={-Latex, thick},
]
  \node[box, fill=red!20] (extract) {1. Extract\\\scriptsize Vendor-Adapter\\\scriptsize (FRED/ECB/Yahoo)};
  \node[box, fill=yellow!20, right=of extract] (validate) {2. Validate Raw\\\scriptsize Pandera Schemas};
  \node[box, fill=orange!20, right=of validate] (transform) {3. Transform\\\scriptsize Polars Mapping\\\scriptsize to canonical schema};
  \node[box, fill=cyan!20, right=of transform] (validate2) {4. Validate Clean\\\scriptsize PK, OHLC, ranges};
  \node[box, fill=green!20, below=of validate] (upsert) {5. Upsert\\\scriptsize SQLAlchemy merge()};
  \node[box, fill=blue!20, right=of upsert] (db) {6. Persist\\\scriptsize PostgreSQL};

  \draw[arrow] (extract) -- (validate);
  \draw[arrow] (validate) -- (transform);
  \draw[arrow] (transform) -- (validate2);
  \draw[arrow] (validate2) -- (upsert);
  \draw[arrow] (upsert) -- (db);
\end{tikzpicture}
\caption{Der ETL-Flow: Sechs Schritte, jeweils mit klarer Verantwortung.
  Im aktuellen Projekt sind die Validate-Schritte vorbereitet, aber
  noch nicht vollständig implementiert (geplant für eine DQ-Iteration).}
\label{fig:etl-flow}
\end{figure}

\section{Der Seeder-Pattern}

Listing~\ref{lst:seeder} zeigt das Standard-Pattern zum
\emph{Seed} (Initial-Befüllung) einer Tabelle. Es wird für
\code{vendors}, \code{instruments} und \code{prices\_daily}
verwendet.

\begin{lstlisting}[language=python, caption={Seeder-Pattern: lookup-or-create},label=lst:seeder]
def upsert_instrument(asset_class, name, currency, symbol):
    """Idempotent: creates the instrument if not in DB, else returns id."""
    with session_scope() as s:
        existing = (
            s.query(Instrument)
            .join(InstrumentIdentifier,
                  InstrumentIdentifier.instrument_id == Instrument.instrument_id)
            .filter(
                InstrumentIdentifier.scheme == "YAHOO",
                InstrumentIdentifier.value == symbol,
            )
            .one_or_none()
        )
        if existing is None:
            inst = Instrument(asset_class=asset_class, name=name, currency=currency)
            s.add(inst)
            s.flush()  # gets us inst.instrument_id
            s.add(InstrumentIdentifier(
                instrument_id=inst.instrument_id,
                scheme="YAHOO",
                value=symbol,
            ))
            s.flush()
            return inst.instrument_id
        return existing.instrument_id
\end{lstlisting}

Die Schlüsseleigenschaft ist \emph{Idempotenz}: Wenn der Seeder
zweimal läuft, werden keine Duplikate angelegt. \code{merge()} allein
reicht dafür nicht aus -- wenn das Objekt keine \code{primary key}
hat, versucht SQLAlchemy ein \code{INSERT} statt eines \code{UPDATE}.
Die explizite \code{SELECT}-Suche im Voraus ist robuster.

\section{Das \code{query\_db.py}-Skript}

Das \code{query\_db.py}-Skript ist die zentrale CLI für Daten-Operationen.
Es kann:

\begin{itemize}
  \item die Datenbank mit Demo-Daten füllen (falls leer),
  \item echte Daten von FRED, ECB und Yahoo holen,
  \item die Daten in tabellarischer Form im Terminal anzeigen.
\end{itemize}

Listing~\ref{lst:query-db} zeigt einen vereinfachten Auszug.

\begin{lstlisting}[language=python, caption={Auszug aus \code{scripts/query\_db.py}},label=lst:query-db]
def main() -> int:
    settings = get_settings()
    configure_logging(settings.log_level)

    with session_scope() as s:
        n_obs = s.execute(select(func.count()).select_from(MacroObservation)).scalar() or 0

    if n_obs == 0 and not args.no_seed:
        # No data: try ECB + FRED in sequence
        if settings.fred_api_key:
            print("seeding real FRED data...")
            _seed_from_fred(settings.fred_api_key)
        print("seeding real ECB data (no API key required)...")
        _seed_from_ecb()
        print("seeding real Yahoo Finance data (AAPL, MSFT, ^GSPC)...")
        _seed_from_yahoo()

    # Display all series with their latest values
    series_list = _fetch_series(args.series)
    _print_summary(series_list, tail=args.tail)
    return 0
\end{lstlisting}

Beachten Sie das \emph{Reihenfolge}-Pattern:

\begin{enumerate}
  \item Prüfe, ob die DB bereits Daten hat (\code{n\_obs > 0}).
  \item Falls leer: Versuche ECB (kein Key) und FRED (falls Key
        gesetzt). Fällt auf Demo-Daten zurück, wenn beide fehlschlagen.
  \item Lese alle Serien und zeige sie formatiert an.
\end{enumerate}

\section{Bulk-Seeding mit \code{seed\_paper\_data.py}}

Für die AEGIS-Paper-Replikation (Kapitel~\ref{ch:analytics}) haben
wir ein eigenes Skript gebaut, das 36~Symbole (31~Aktien + 5~Benchmarks)
in einem einzigen Aufruf von \code{yfinance} lädt.

\begin{lstlisting}[language=shell, caption={Bulk-Seed-Skript aufrufen}]
# Einmaliger Aufruf
.venv/bin/python scripts/seed_paper_data.py

# Mit benutzerdefiniertem Zeitraum
.venv/bin/python scripts/seed_paper_data.py --start 2006-01-01

# Ohne Benchmarks
.venv/bin/python scripts/seed_paper_data.py --no-benchmarks
\end{lstlisting}

Das Skript nutzt die \code{fetch\_many}-Bulk-Methode des
Yahoo-Vendors (Listing~\ref{lst:yahoo-fetch}) -- 36~Symbole
werden in einem einzigen \code{yf.download}-Aufruf geholt, was
~12~Sekunden statt 36~einzelner Calls benötigt. Insgesamt
werden 185.220~Zeilen in die \code{prices\_daily}-Tabelle
geschrieben.

\section{DQ-Checks (Vorbereitung)}

Das Projekt hat Skelette für Data-Quality-Checks, die in einer
späteren Iteration implementiert werden sollen. Die
DQ-Tabellen (\code{dq\_runs}, \code{dq\_findings}) existieren bereits.

Folgende Checks sind geplant:

\begin{itemize}
  \item \textbf{PK-Eindeutigkeit}: \code{(instrument\_id, date)}
        darf nicht doppelt vorkommen.
  \item \textbf{OHLC-Konsistenz}: \code{low} $\leq$ \code{open}, \code{close}
        $\leq$ \code{high} (keine \emph{inverted candles}).
  \item \textbf{Keine negativen Preise/Volumes}.
  \item \textbf{Adjusted-Close-Konsistenz}: Wenn \code{corporate\_actions}
        verfügbar ist, muss \code{adjusted\_close} das Produkt aus
        \code{close} und allen Splits/Dividenden sein.
  \item \textbf{Lücken-Detection}: Mehr als 5~aufeinanderfolgende
        Handelstage ohne Eintrag -- ein \emph{DQ-Finding}.
  \item \textbf{Ausreißer-Detection}: Tagesreturn $> 6\sigma$
        gegenüber rollendem 252-Tage-Fenster -- \emph{flag}, nicht
        löschen.
  \item \textbf{Currency-Coverage}: Jedes \code{instrument} hat
        \code{currency} gesetzt.
  \item \textbf{Vendor-Coverage}: Jede Zeile in \code{prices\_daily}
        hat \code{vendor\_id} gesetzt.
\end{itemize}

\section{Idempotenz}

\begin{merke}
Alle Seeder sind \emph{idempotent}. Sie können beliebig oft
ausgeführt werden, ohne Duplikate zu erzeugen. Dies ist entscheidend
für die Inkrementelle-Aktualisierung: ein täglicher Cron-Job ruft
den Seeder auf, und nur neue Daten werden tatsächlich geschrieben.
\end{merke}

Listing~\ref{lst:idempotency} zeigt das zentrale Idempotenz-Pattern.

\begin{lstlisting}[language=python, caption={Idempotenz-Pattern: SELECT vor INSERT},label=lst:idempotency]
def upsert_prices(instrument_id, vendor_id, df) -> int:
    if df.is_empty():
        return 0
    rows = 0
    with session_scope() as s:
        for r in df.iter_rows(named=True):
            s.merge(PricesDaily(
                instrument_id=instrument_id,
                date=r["date"],
                open=r.get("open"),
                high=r.get("high"),
                low=r.get("low"),
                close=r["close"],
                adjusted_close=r.get("adjusted_close"),
                volume=r.get("volume"),
                vendor_id=vendor_id,
            ))
            rows += 1
    return rows
\end{lstlisting}

\section{Failure-Isolation}

\begin{definition}
\textbf{Failure-Isolation}
\emph{Failure-Isolation} bedeutet, dass der Fehler eines Vendors
\emph{nicht} den gesamten Pipeline-Lauf abbricht. Wenn Yahoo
gerade ein Problem hat, sollen ECB und FRED trotzdem seeden.
\end{definition}

Das Skript \code{query\_db.py} setzt das um, indem jeder
Vendor-Seeder in einem separaten \code{session\_scope} läuft. Ein
Fehler in einem Vendor wird geloggt, der nächste Vendor läuft trotzdem.

\section{Was als Nächstes?}

Die ETL-Pipeline produziert Daten in \code{prices\_daily} und
\code{macro\_observations}. Im nächsten Kapitel sehen wir, wie diese
Daten in einem \emph{Jupyter-Notebook} visualisiert werden können.

% =====================================================================
% Kapitel 7: Notebook-Generator
% =====================================================================

\chapter{Notebook-Generator}

\label{ch:notebook}

Das \emph{Jupyter-Notebook} ist die komfortabelste Form, um
Finanzmarktdaten zu visualisieren. Es ermöglicht interaktive
Charts, schrittweises Erforschen und einfaches Teilen der Ergebnisse
als HTML.

\section{Überblick: Pipeline}

\begin{figure}[ht]
\centering
\begin{tikzpicture}[
  node distance=0.8cm and 1.2cm,
  box/.style={rectangle, draw, thick, rounded corners, minimum width=2.5cm, minimum height=0.8cm, align=center, font=\small},
  file/.style={ellipse, draw, thick, minimum width=2.5cm, font=\small, align=center},
  arrow/.style={-Latex, thick},
]
  \node[file, fill=cyan!20] (build) {build\_notebook.py};
  \node[file, fill=yellow!20, right=of build] (ipynb) {\code{01\_macro\_overview.ipynb}};
  \node[file, fill=green!20, right=of ipynb] (html) {\code{01\_macro\_overview.html}};

  \draw[arrow] (build) -- node[above, font=\scriptsize] {\code{nbformat}} (ipynb);
  \draw[arrow] (ipynb) -- node[above, font=\scriptsize] {\code{nbconvert --execute}} (html);
\end{tikzpicture}
\caption{Build-Pipeline für das Notebook. \code{build\_notebook.py}
  erzeugt die \code{.ipynb} programmatisch aus Cell-Definitionen.
  \code{jupyter nbconvert} führt die Cells aus und exportiert HTML.}
\label{fig:notebook-pipeline}
\end{figure}

\section{Warum programmatisch generieren?}

\begin{warnung}
\rmfamily\itshape
\foreign{Notebook-Ausführungs-Ausgaben veralten schnell}. Wenn ein
Notebook committet wird, sind die eingebetteten Plots und Tabellen
\emph{für immer} eingefroren -- egal, wie sich die Daten dahinter
ändern. Das führt zu \emph{stale-notebooks}, in denen der Leser
nicht erkennt, dass die Daten veraltet sind.
\end{warnung}

Die Lösung: \emph{Trennung} zwischen \code{Quellcode} (Cell-Definitionen
in Python) und \emph{Output} (gerendertes Notebook + HTML). Die
\code{.ipynb}-Datei ist ein \emph{Build-Artefakt}, das jederzeit
regeneriert werden kann.

\section{Der Build-Prozess}

Listing~\ref{lst:build} zeigt das Herzstück von
\code{scripts/build\_notebook.py}.

\begin{lstlisting}[language=python, caption={Notebook programmatisch erzeugen},label=lst:build]
import nbformat as nbf

CELLS = [
    nbf.v4.new_markdown_cell("# deephold_db — Macro Overview"),
    nbf.v4.new_code_cell("from deephold_db.db import session_scope, MacroObservation"),
    nbf.v4.new_code_cell("..."),
    nbf.v4.new_markdown_cell("## Plots"),
    nbf.v4.new_code_cell("fig = make_subplots(...)"),
    nbf.v4.new_code_cell("fig.show()"),
    nbf.v4.new_markdown_cell("## Summary"),
    nbf.v4.new_code_cell("summary = pl.DataFrame(rows); print(summary)"),
]

def main():
    nb = nbf.v4.new_notebook()
    nb.cells = CELLS
    nb.metadata = {"kernelspec": {"name": "python3"}}
    nbf.write(nb, NOTEBOOK_PATH)
    subprocess.run(["jupyter", "nbconvert", "--execute", "--inplace", str(NOTEBOOK_PATH)])
    subprocess.run(["jupyter", "nbconvert", "--to", "html", str(NOTEBOOK_PATH)])
\end{lstlisting}

Der Build-Prozess ist in vier Schritte:

\begin{enumerate}
  \item \textbf{Cell-Definition}: Python-Liste von
        \code{nbf.v4.new\_markdown\_cell()} und
        \code{nbf.v4.new\_code\_cell()} Objekten.
  \item \textbf{Notebook schreiben}: \code{nbf.write()} erzeugt
        eine valides \code{.ipynb}-JSON-Datei.
  \item \textbf{Execute}: \code{jupyter nbconvert --execute}
        führt alle Code-Cells aus und bettet die Outputs in die
        \code{.ipynb} ein.
  \item \textbf{Export}: \code{jupyter nbconvert --to html}
        erzeugt eine statische HTML-Datei, die ohne Jupyter
        geöffnet werden kann.
\end{enumerate}

\section{Plotly 3$\times$2-Grid}

Das Notebook zeigt 6~Subplots in einem 3$\times$2-Grid
(Abbildung~\ref{fig:notebook-grid}):

\begin{figure}[ht]
\centering
\begin{tikzpicture}[
  box/.style={rectangle, draw, thick, minimum width=4cm, minimum height=1.5cm, align=center, font=\small},
  arrow/.style={-Latex, thick},
]
  \node[box] (y) at (0, 2) {Yields\\(DGS3MO, DGS10, FEDFUNDS)};
  \node[box] (i) at (5, 2) {Inflation\\(CPI vs HICP YoY)};
  \node[box] (fx) at (0, 0) {FX Major\\(USD/EUR, GBP/EUR, JPY/EUR)};
  \node[box] (un) at (5, 0) {Unemployment\\(UNRATE)};
  \node[box] (ep) at (0, -2) {Equity Prices\\(AAPL, MSFT, S\&P 500)};
  \node[box] (ev) at (5, -2) {Equity Volume\\(AAPL, MSFT)};
\end{tikzpicture}
\caption{3$\times$2 Plotly-Subplot-Grid im Notebook. Die obere
  Hälfte zeigt Makro-Daten, die untere Hälfte Equity-Daten.}
\label{fig:notebook-grid}
\end{figure}

Jeder Subplot hat seine eigenen \emph{Y-Achsen-Skalen}, weil die
Größenordnungen fundamental verschieden sind (Zinssätze in
\,\%, Aktienkurse in USD, Volumen in Milliarden).

\section{HTML-Export}

Der HTML-Export hat zwei Verwendungszwecke:

\begin{itemize}
  \item \textbf{Sharing}: Man kann die HTML-Datei per Mail
        verschicken oder auf einem Webserver ablegen, ohne dass
        der Empfänger Jupyter installieren muss.
  \item \textbf{Caching}: Die HTML-Datei ist ein \emph{Snapshot}
        der Daten zum Build-Zeitpunkt. Wenn die DB später
        aktualisiert wird, bleibt der HTML-Snapshot konsistent.
\end{itemize}

\begin{merke}
Die HTML-Datei \emph{referenziert} plotly.js von einem CDN. Ohne
Internet-Zugang beim Öffnen der HTML-Datei sind die Plots
\emph{nicht} sichtbar -- die zugrundeliegenden Daten sind aber
weiterhin im HTML-Markup eingebettet.
\end{merke}

\section{Was als Nächstes?}

Das Notebook ist eine schöne \emph{Passive-View} der Daten. Für
\emph{aktive Exploration} -- Springen zwischen Serien, Suche,
Keyboard-Navigation -- brauchen wir etwas anderes: das
\emph{OpenTUI-Explorer}, das im nächsten Kapitel vorgestellt wird.

% =====================================================================
% Kapitel 8: OpenTUI-Explorer
% =====================================================================

\chapter{OpenTUI-Explorer}

\label{ch:tui}

Der \emph{OpenTUI-Explorer} ist die schnellste und interaktivste
Möglichkeit, die gesammelten Daten zu durchsuchen. Er läuft direkt
im Terminal, ohne Jupyter, ohne Notebook, ohne Python-Overhead im
Hot-Path.

\section{Was ist OpenTUI?}

\begin{definition}
\textbf{OpenTUI}
OpenTUI (https://opentui.com) ist ein von Anomaly (https://anoma.ly)
entwickeltes TUI-Framework. Es besteht aus:

\begin{itemize}
  \item einem in \emph{Zig} geschriebenen \emph{Core}
        (schnell, plattformübergreifend, nativ),
  \item einer \emph{C-ABI} für Bindings aus jeder Sprache,
  \item \emph{TypeScript}-Bindings mit React- und Solid.js-Support.
\end{itemize}

Der Kern nutzt \emph{Yoga} (Facebooks Flexbox-Engine) für Layouts
und rendert das Ergebnis direkt in einen Terminal-Buffer.
\end{definition}

\begin{merke}
OpenTUI ist \emph{die gleiche Engine}, die \foreign{opencode}
(powered by Anomaly) und \foreign{terminal.shop} antreibt. Es
befindet sich in aktiver Entwicklung, kann also Breaking Changes
enthalten.
\end{merke}

\section{Warum TUI statt GUI?}

\begin{figure}[ht]
\centering
\begin{tabular}{lll}
  \toprule
  \textbf{Aspekt} & \textbf{TUI} & \textbf{GUI (z.\,B.\ Web-App)} \\
  \midrule
  Startup-Zeit          & $<$100 ms (Bun)         & $>$1 s (Browser) \\
  Speicherverbrauch     & $\sim$50 MB               & $\sim$300 MB \\
  Tastatur-Effizienz    & \foreign{Vim-style}      & Maus \\
  SSH-/Server-tauglich  & ja                       & nur mit X-Forwarding \\
  Visueller Reiz        & niedrig                  & hoch \\
  Lernkurve             & hoch (Tasten lernen)     & niedrig (klicken) \\
  \bottomrule
\end{tabular}
\caption{TUI vs.\ GUI}
\label{fig:tui-vs-gui}
\end{figure}

Für die tägliche Exploration von Finanzmarktdaten -- springen
zwischen 12~Serien, monatliche Updates prüfen, Daten vergleichen --
ist eine TUI das richtige Werkzeug.

\section{Architektur des TUI-Layers}

Das TUI ist ein \emph{eigenständiger TypeScript-Subtree} im
Projekt-Root:

\begin{lstlisting}[language=shell, caption={TUI-Verzeichnisstruktur}]
tui/
├── package.json            # Bun + React 19 + OpenTUI 0.4
├── tsconfig.json
├── bun.lock
├── src/
│   ├── index.tsx          # Entry: createCliRenderer + createRoot
│   ├── App.tsx            # Root-Komponente (Layout, State, Keyboard)
│   ├── components/
│   │   ├── SeriesList.tsx   # Linke Seite: 12 Serien
│   │   └── SeriesDetail.tsx # Rechte Seite: Stats + Tail
│   ├── data/
│   │   ├── db.ts          # pg-Pool
│   │   └── queries.ts     # 4 SQL-Queries
│   └── types.ts           # TypeScript-Typen
├── tests/                 # bun test (Unit + Integration)
└── scripts/check-db.ts    # DB-Connectivity-Test
\end{lstlisting}

\section{Die Komponenten}

\subsection{Entry Point: \code{index.tsx}}

\begin{lstlisting}[language=typescript, caption={TUI-Entry-Point},label=lst:tui-entry]
import { createCliRenderer } from "@opentui/core";
import { createRoot } from "@opentui/react";
import { App } from "./App";
import { closePool } from "./data/db";

const renderer = await createCliRenderer({ exitOnCtrlC: true });
const root = createRoot(renderer);
root.render(<App />);
\end{lstlisting}

Der Entry-Point ist minimal:

\begin{enumerate}
  \item \code{createCliRenderer()} -- startet den OpenTUI-Renderer,
        der den Terminal-Buffer verwaltet.
  \item \code{createRoot(renderer)} -- erzeugt einen React-Root,
        der direkt in den OpenTUI-Buffer rendert.
  \item \code{root.render(<App />)} -- mountet die App-Komponente.
\end{enumerate}

\subsection{Root: \code{App.tsx}}

Die Root-Komponente verwaltet drei Dinge: \emph{State} (welche Serien
sind in der DB, welche ist selektiert), \emph{Keyboard-Handling}
und das \emph{Layout} (2-Spalten).

\begin{lstlisting}[language=typescript, caption={TUI-App: State + Keyboard},label=lst:tui-app]
function App() {
  const renderer = useRenderer();
  const [series, setSeries] = useState<SeriesRow[]>([]);
  const [selectedIndex, setSelectedIndex] = useState(0);

  useEffect(() => {
    listSeries().then(setSeries).catch(/* ... */);
  }, []);

  useKeyboard((key) => {
    if (key.name === "q") renderer.destroy();
    if (key.name === "r") setRefreshToken(t => t + 1);
    if (key.name === "up" || key.name === "k") setSelectedIndex(i => Math.max(0, i - 1));
    if (key.name === "down" || key.name === "j") setSelectedIndex(i => Math.min(series.length - 1, i + 1));
  });

  return (
    <box flexDirection="row">
      <SeriesList series={series} selectedIndex={selectedIndex} />
      {series[selectedIndex] && <SeriesDetail series={series[selectedIndex]} />}
    </box>
  );
}
\end{lstlisting}

\subsection{SeriesList -- Linke Spalte}

Die \code{SeriesList}-Komponente rendert alle geladenen Serien als
\emph{scrollbare} Liste. Jede Zeile enthält:

\begin{itemize}
  \item \textbf{Badge}: \code{M} für Macro (FRED/ECB), \code{E} für
        Equity (Yahoo).
  \item \textbf{Series-ID}: z.\,B.\ \code{FRED:DGS3MO} oder
        \code{ECB:EXR:USD.EUR.SP00.A}.
  \item \textbf{Vendor-Tag}: \code{fred}, \code{ecb}, \code{yahoo}.
  \item \textbf{Row-Count}: rechtsbündig für schnelles Scannen.
\end{itemize}

\subsection{SeriesDetail -- Rechte Spalte}

Die \code{SeriesDetail}-Komponente zeigt für die ausgewählte Serie:

\begin{itemize}
  \item \textbf{Header}: Series-ID + Name.
  \item \textbf{Stats-Panel}: rows, first/last date, unit, source,
        min/max/mean.
  \item \textbf{Tail-Tabelle}: die letzten 30~Werte, bei Macro
        als \code{date, value}, bei Equity als
        \code{date, close, adj\_close, vol}.
\end{itemize}

\section{Die SQL-Queries}

Das TUI spricht \emph{direkt} mit der PostgreSQL-DB via
\code{node-postgres} -- es gibt keinen Python-Overhead im Hot-Path.

\begin{lstlisting}[language=typescript, caption={4 SQL-Queries im TUI}]
const LIST_SERIES_SQL = `
  SELECT 'macro' AS kind, series_id AS id, name, source, frequency, unit, COUNT(*)::int AS n
  FROM macro_observations mo
  JOIN macro_series ms USING (series_id)
  GROUP BY series_id, name, source, frequency, unit
  UNION ALL
  SELECT 'equity' AS kind, COALESCE(ii.value, i.instrument_id::text) AS id,
         i.name, 'yahoo' AS source, 'D' AS frequency, i.currency AS unit, COUNT(*)::int AS n
  FROM prices_daily pd
  JOIN instruments i USING (instrument_id)
  LEFT JOIN instrument_identifiers ii
    ON ii.instrument_id = i.instrument_id AND ii.scheme = 'YAHOO'
  GROUP BY i.instrument_id, ii.value, i.name, i.currency
  ORDER BY 1, 2
`;

const MACRO_TAIL_SQL = `
  SELECT to_char(date, 'YYYY-MM-DD') AS date, value
  FROM macro_observations WHERE series_id = $1
  ORDER BY date DESC LIMIT $2
`;

const EQUITY_TAIL_SQL = `
  SELECT to_char(pd.date, 'YYYY-MM-DD') AS date, pd.close, pd.adjusted_close, pd.volume
  FROM prices_daily pd
  JOIN instrument_identifiers ii
    ON ii.instrument_id = pd.instrument_id AND ii.scheme = 'YAHOO'
  WHERE ii.value = $1
  ORDER BY pd.date DESC LIMIT $2
`;

const MACRO_STATS_SQL = `
  SELECT COUNT(*)::int AS n, to_char(MIN(date), 'YYYY-MM-DD') AS first,
         to_char(MAX(date), 'YYYY-MM-DD') AS last, MIN(value) AS min, MAX(value) AS max, AVG(value) AS mean
  FROM macro_observations WHERE series_id = $1
`;
\end{lstlisting}

Beachten Sie:

\begin{itemize}
  \item \textbf{UNION ALL} -- eine einzige Query listet Macro- und
        Equity-Serien in einer Liste.
  \item \code{COALESCE(ii.value, i.instrument\_id::text)} -- für
        Equity-Serien wird der YAHOO-Ticker (\code{AAPL}) statt
        der numerischen \code{instrument\_id} (\code{1}) angezeigt.
  \item \code{LEFT JOIN instrument\_identifiers} -- damit auch
        Equity-Serien ohne YAHOO-Identifier angezeigt würden (was
        aktuell nicht vorkommt, aber defensiv ist).
\end{itemize}

\section{Tastatur-Handling}

\begin{table}[ht]
\centering
\small
\begin{tabular}{ll}
  \toprule
  \textbf{Taste} & \textbf{Aktion} \\
  \midrule
  \code{$\uparrow$} / \code{k}  & Selection nach oben \\
  \code{$\downarrow$} / \code{j} & Selection nach unten \\
  \code{g} / \code{Home}  & Zum Anfang \\
  \code{G} / \code{End}   & Zum Ende \\
  \code{r}                & Re-Query (für Live-Refresh) \\
  \code{q} / \code{Ctrl-C}  & Quit \\
  \bottomrule
\end{tabular}
\caption{Tastatur-Bindings des TUI. Die \code{j}/\code{k}-Variante
  ist Vim-style für Entwickler, die diese Tasten gewohnt sind.}
\label{tab:tui-keys}
\end{table}

\begin{lstlisting}[language=typescript, caption={Keyboard-Handler},label=lst:tui-keys]
useKeyboard((key) => {
  if (key.name === "q" || (key.ctrl && key.name === "c")) {
    renderer.destroy();
    setTimeout(() => process.exit(0), 50);
    return;
  }
  if (key.name === "r") {
    setRefreshToken(t => t + 1);
    return;
  }
  if (key.name === "up" || key.name === "k") {
    setSelectedIndex(i => Math.max(0, i - 1));
    return;
  }
  // ... etc.
});
\end{lstlisting}

\section{Tests}

Das TUI hat 27~Tests in 4~Dateien:

\begin{itemize}
  \item \code{src/data/queries.test.ts} (9~Tests) -- SQL-Strings,
        Row-Mapping, Param-Übergabe.
  \item \code{src/data/db.test.ts} (3~Tests) -- Public-API des
        \code{db}-Moduls.
  \item \code{src/types/format.test.ts} (8~Tests) -- Pure-Funktionen
        (\code{fmtNumber}, \code{fmtPct}).
  \item \code{tests/integration.test.ts} (7~Tests) -- Echtes
        Postgres, 12~Series, OHLCV-Validierung.
\end{itemize}

\begin{lstlisting}[language=bash, caption={TUI-Tests laufen}]
cd tui && bun test
# 27 pass, 0 fail, 203 expect(), 4 files
\end{lstlisting}

\section{Was als Nächstes?}

Die TUI ist die \emph{operative Oberfläche} des Projekts. Das
nächste Kapitel zeigt, wie wir die Datenqualität \emph{testen} --
nicht nur funktional (geht der Code?), sondern auch inhaltlich
(stimmen die Daten?).

% =====================================================================
% Kapitel 9: Tests
% =====================================================================

\chapter{Tests}

\label{ch:tests}

Tests sind das Sicherheitsnetz des Projekts. Sie prüfen, ob der
Code das tut, was er tun soll -- und verhindern, dass spätere
Änderungen bestehende Funktionalität brechen.

\section{Test-Philosophie}

\begin{merke}
Das Projekt folgt der \emph{Pyramide} der Tests:

\begin{itemize}
  \item \textbf{Viele} schnelle \emph{Unit-Tests} (einzelne Funktionen).
  \item \textbf{Wenige} \emph{Integration-Tests} (mehrere
        Komponenten zusammen).
  \item \textbf{Manuelle} \emph{Checks} für Dinge, die schwer zu
        automatisieren sind (visuelle TUI-Korrektheit).
\end{itemize}
\end{merke}

\section{Python-Tests mit pytest}

\subsection{Überblick}

\begin{lstlisting}[language=bash, caption={Python-Tests laufen}]
.venv/bin/pytest
# 55 passed
\end{lstlisting}

55~Tests sind in der Suite (Stand: Handbuch-Erstellung). Sie sind
in folgende Module aufgeteilt:

\begin{itemize}
  \item \code{tests/test\_config.py} -- Konfigurationsladen.
  \item \code{tests/test\_models.py} -- SQLAlchemy-ORM-Registrierung.
  \item \code{tests/test\_fred.py} -- FRED-Vendor (mit gemocktem
        HTTP).
  \item \code{tests/test\_ecb.py} -- ECB-Vendor (mit gemocktem HTTP).
  \item \code{tests/test\_yahoo.py} -- Yahoo-Vendor (mit
        gemocktem yfinance).
  \item \code{tests/analytics/test\_metrics.py} -- Returns + Metriken.
\end{itemize}

\subsection{Test-Fixtures}

\begin{definition}
\textbf{pytest Fixture}
Eine \emph{Fixture} ist eine wiederverwendbare Test-Setup-Funktion.
Sie wird mit \code{@pytest.fixture} dekoriert und kann in Tests
\emph{injiziert} werden, indem sie als Argument mit dem gleichen
Namen wie die Fixture-Funktion deklariert wird.
\end{definition}

\begin{lstlisting}[language=python, caption={pytest-Fixture für Settings-Cache-Reset},label=lst:fixture]
@pytest.fixture(autouse=True)
def _reset_settings_cache():
    """Reset the pydantic-settings lru_cache around each test."""
    get_settings.cache_clear()
    yield
    get_settings.cache_clear()
\end{lstlisting}

Diese \code{autouse=True}-Fixture wird \emph{vor jedem Test}
automatisch ausgeführt und stellt sicher, dass der
\code{lru\_cache} der \code{get\_settings()}-Funktion geleert
ist -- sonst würden Test-Pollution und unstabile Ergebnisse
entstehen.

\subsection{Mocking}

\begin{definition}
\textbf{Mock}
Ein \emph{Mock} ist ein gefälschtes Objekt, das das Verhalten eines
realen Objekts nachahmt. In Tests verwenden wir Mocks, um externe
Abhängigkeiten (HTTP-Server, Datenbanken) zu simulieren, ohne
\emph{tatsächlich} mit ihnen zu sprechen -- das macht Tests
\emph{ schnell}, \emph{ deterministisch } und \emph{ offline
  lauffähig}.
\end{definition}

Listing~\ref{lst:mock-fred} zeigt einen typischen Mock-Test für
den FRED-Vendor.

\begin{lstlisting}[language=python, caption={FRED-Vendor mit gemocktem HTTP},label=lst:mock-fred]
def test_fred_fetch_parses_response(fake_fred_observations, monkeypatch):
    monkeypatch.setenv("FRED_API_KEY", "test-key")
    v = FredVendor()
    v._get = MagicMock(return_value=fake_fred_observations)
    df = v.fetch("DGS3MO", date(2024, 6, 1), date(2024, 6, 30))
    assert df.height == 3
    assert df["value"][0] == pytest.approx(5.02)
    v._get.assert_called_once()
\end{lstlisting}

Die Strategie: \code{v.\_get} wird durch einen \code{MagicMock}
ersetzt, der \emph{jeden} Aufruf mit dem vorbereiteten Fake-Response
beantwortet. So können wir die \emph{Response-Parsing}-Logik
isoliert testen, ohne einen echten HTTP-Server zu kontaktieren.

\section{TypeScript-Tests mit bun test}

\subsection{Überblick}

\begin{lstlisting}[language=bash, caption={TUI-Tests laufen}]
cd tui && bun test
# 27 pass, 0 fail
\end{lstlisting}

\subsection{Mocking in bun test}

Bun's Test-Runner hat eine Jest-kompatible API:
\code{describe}, \code{test}, \code{expect}, \code{spyOn}, \code{mock}.

\begin{lstlisting}[language=typescript, caption={Mocking der pg-Layer im TUI},label=lst:bun-mock]
import { describe, it, expect, spyOn, type Mock } from "bun:test";
import * as dbModule from "./db";
import { listSeries } from "./queries";

let spies: Mock<typeof dbModule.query>[] = [];

function mockQuery<T>(rows: T[]): Mock<typeof dbModule.query> {
  const spy = spyOn(dbModule, "query") as Mock<typeof dbModule.query>;
  spy.mockReset();
  spy.mockResolvedValue(rows as never);
  spies.push(spy);
  return spy;
}

afterEach(() => {
  for (const s of spies) s.mockRestore();
  spies = [];
});

describe("listSeries", () => {
  it("calls query() and returns the rows", async () => {
    const rows: SeriesRow[] = [{...}];
    const spy = mockQuery(rows);
    const out = await listSeries();
    expect(out).toEqual(rows);
    expect(spy).toHaveBeenCalledTimes(1);
    const sql = spy.mock.calls[0]?.[0] as string;
    expect(sql).toContain("UNION ALL");
  });
});
\end{lstlisting}

Wichtige Pattern:

\begin{itemize}
  \item \code{spyOn} ersetzt eine Methode durch einen Spy, der
        Aufrufe aufzeichnet.
  \item \code{mockResolvedValue} setzt den Rückgabewert.
  \item \code{afterEach} restauriert die Spies, um Test-Pollution zu
        vermeiden.
\end{itemize}

\subsection{Integration-Tests}

\begin{lstlisting}[language=typescript, caption={Integration-Test mit echtem Postgres}]
beforeAll(async () => {
  try {
    await query("SELECT 1::int AS x");
    dbReachable = true;
  } catch (e) {
    console.warn("Postgres unreachable, tests will no-op.");
  }
});

function guarded(name, fn) {
  return async () => {
    if (!dbReachable) {
      console.warn(`skip ${name}: DB unreachable`);
      return;
    }
    await fn();
  };
}

it("equity rows carry the YAHOO ticker (AAPL/MSFT/^GSPC)",
   guarded("equity-tickers", async () => {
     const rows = await listSeries();
     const tickers = rows
       .filter(r => r.kind === "equity")
       .map(r => r.id)
       .sort();
     expect(tickers).toEqual(["AAPL", "MSFT", "^GSPC"]);
   }));
\end{lstlisting}

Integration-Tests prüfen das Verhalten gegen die echte
PostgreSQL-DB. Sie sind \emph{no-op} (nicht \emph{fail}), wenn
die DB nicht erreichbar ist -- so bleibt die CI grün, auch wenn
Postgres temporär down ist.

\section{Manuelle Test-Checklist}

Manche Dinge sind schwer zu automatisieren -- visuelle Korrektheit
einer TUI, Tastatur-Feel, Verhalten unter ungewöhnlichen
Bedingungen. Dafür gibt es die manuelle Checkliste in
\code{tui/README.md}.

\begin{beispiel}
\textbf{Auszug aus der manuellen TUI-Checklist}
  \begin{enumerate}
    \item \code{make tui} starten
    \item Header zeigt "12 series" (nicht "loading...")
    \item Erste Zeile (\code{AAPL}) ist selektiert ($\triangleright$)
    \item \code{$\downarrow$} 4x drücken: Selection wandert
    \item \code{r} drücken: Series-Liste wird neu geladen
    \item \code{q} drücken: TUI beendet sauber
  \end{enumerate}
\end{beispiel}

\section{Code-Qualität}

\begin{definition}
\textbf{Linter}
Ein \emph{Linter} ist ein statisches Code-Analyse-Tool, das nach
\emph{Potentiellen Fehlern}, \emph{Stil-Verstößen} und
\emph{Sicherheitsproblemen} sucht, ohne den Code auszuführen.
\end{definition}

Dieses Projekt verwendet:

\begin{itemize}
  \item \textbf{ruff} -- ein extrem schneller Python-Linter
        (in Rust geschrieben). Ersetzt \code{flake8}, \code{isort},
        \code{black} und Teile von \code{pylint}.
  \item \textbf{tsc} -- der TypeScript-Compiler im
        \code{--noEmit}-Modus (nur Typ-Check, keine Ausgabe).
\end{itemize}

\begin{lstlisting}[language=bash, caption={Linter ausführen}]
# Python
.venv/bin/ruff check scripts/ src/ tests/
.venv/bin/ruff format --check scripts/ src/ tests/

# TypeScript
cd tui && bun run typecheck
\end{lstlisting}

\begin{warnung}
Beachte: Das Projekt verwendet \emph{keine} Typen-Annotationen für
\`mypy\` -- die Typen sind in \code{pyproject.toml} unter
\code{[tool.mypy]} deaktiviert. Dies ist eine bewusste
Vereinfachung für ein Retail-Projekt; für ein
\emph{Produktions}-System würde ich strikte Typen empfehlen.
\end{warnung}

\section{Was als Nächstes?}

Mit Tests und Linting als Sicherheitsnetz sind wir bereit für das
\emph{letzte und ambitionierteste} Stück: die
AEGIS-Paper-Replikation (Kapitel~\ref{ch:analytics}). Dort
verbinden wir alle bisherigen Komponenten -- Vendor-Adapter,
ETL-Pipeline, Datenmodell -- zu einer kohärenten quantitativen
Analyse.

% =====================================================================
% Kapitel 10: Analytics und AEGIS-Paper-Replikation
% =====================================================================

\chapter{Analytics und AEGIS-Paper-Replikation}

\label{ch:analytics}

Das letzte und ambitionierteste Stück des Projekts: die
\emph{Replikation} eines quantitativen Forschungs-Papers.
Wir verwenden die gesamte bisherige Infrastruktur, um eine
moderne \emph{Momentum-Strategie} zu implementieren und ihre
Performance gegen reale Benchmarks zu messen.

\section{Das AEGIS-Paper}

\begin{definition}
\textbf{AEGIS}
AEGIS (\emph{Adaptive Equity Generation and Immunisation System})
ist ein Framework für quantitative Aktien-Strategien, vorgestellt
in Chakraborty und Singh (2026), \foreign{Taming the Black Swan:
A Momentum-Gated Hierarchical Optimisation Framework for Asymmetric
Alpha Generation} (arXiv:2604.09060v1). Das Paper
berichtet eine 20-Jahres-Backtest-Performance von
\emph{15{,}41\,\% CAGR} bei \emph{-28{,}89\,\% MaxDD} -- die
NASDAQ-100-Rendite mit deutlich weniger Drawdown.
\end{definition}

\subsection{Die 3-Layer-Architektur}

\begin{figure}[ht]
\centering
\begin{tikzpicture}[
  node distance=0.6cm,
  box/.style={rectangle, draw, thick, rounded corners, minimum width=3.5cm, minimum height=1cm, align=center, font=\small},
  arrow/.style={-Latex, thick},
  layer/.style={box, fill=#1!15},
]
  \node[layer=blue] (sig) {1. Signal Generation\\\scriptsize VAM = R/$\sigma$\\\scriptsize Anchor Triad};
  \node[layer=green, below=of sig] (imm) {2. Immunisation\\\scriptsize Minimax-Correlation\\\scriptsize Diversifier-Selection};
  \node[layer=orange, below=of imm] (alloc) {3. Allocation\\\scriptsize SLSQP\\\scriptsize Sortino-Maximierung};

  \draw[arrow] (sig) -- (imm);
  \draw[arrow] (imm) -- (alloc);
\end{tikzpicture}
\caption{Die 3-Layer-Architektur von AEGIS. Im aktuellen Projekt
  implementieren wir nur Layer~1 (VAM) -- Layer 2 (Minimax) und
  Layer 3 (SLSQP) sind als nächste Iterationen geplant.}
\label{fig:aegis-layers}
\end{figure}

\section{Das VAM-Signal}

\begin{definition}
\textbf{VAM (Volatility-Adjusted Momentum)}
Das VAM-Signal pro Aktie $i$ zum Zeitpunkt $t$:

\[
  \text{VAM}_i(t) = \frac{R_i(t-L:t-s)}{\sigma_i(t-L:t-s)}
\]

wobei $R_i$ die kumulative Log-Rendite und $\sigma_i$ die annualisierte
Volatilität über das Fenster $[t-L, t-s)$ sind. $L$ ist die
Lookback-Periode (Default: 252~Tage = 12~Monate), $s$ der
Skip-Month ($s=21$ Tage = 1~Monat).
\end{definition}

Der Skip-Month ist entscheidend: die jüngsten 21~Tage enthalten
oft \emph{Mean-Reversion}-Rauschen, das das Signal verzerren würde.

\begin{lstlisting}[language=python, caption={VAM-Implementierung},label=lst:vam]
def vam(close, window=252, skip=21):
    r = cumulative_log_return(close, window, skip)
    s = realized_vol(close, window, skip)
    if math.isnan(r) or math.isnan(s) or s == 0:
        return float("nan")
    return r / s
\end{lstlisting}

\section{Walk-Forward-Backtest}

\begin{definition}
\textbf{Walk-Forward-Backtest}
Ein Walk-Forward-Backtest vermeidet \emph{Look-Ahead-Bias} durch
\emph{zeitliche Trennung} von Training und Test:

\begin{enumerate}
  \item Trainiere auf $[t-L, t]$ (z.\,B.\ 12~Monate).
  \item Evaluiere auf $[t, t+1]$ (z.\,B.\ 1~Monat).
  \item Wandere um $t+1$ weiter, wiederhole.
\end{enumerate}

Der Algorithmus hat zu jedem Zeitpunkt nur Zugriff auf Daten, die
\emph{tatsächlich} verfügbar waren -- das eliminiert die
Look-Ahead-Bias vollständig.
\end{definition}

\begin{lstlisting}[language=python, caption={Walk-Forward-Backtest: Top-N by VAM}]
def run_vam_topn_backtest(prices, rebalance_every=21, top_n=10):
    n = min_len = min(len(s) for s in prices.values())
    start_idx = window + skip + 1  # need history

    for end_idx in range(start_idx, n, rebalance_every):
        # 1. Score each stock with VAM
        scores = {sym: vam(_slice_window(p, end_idx, window, skip))
                   for sym, p in prices.items()}

        # 2. Top-N by VAM, with positive gate
        top = [s for s in rank_top_n(scores, top_n*2) if scores[s] > 0][:top_n]

        # 3. Equal-weight, hold for one month
        for i in range(end_idx, min(end_idx + rebalance_every, n)):
            r = sum(log_rets[sym][i] / len(top) for sym in top)
            rets_in_period.append(r)

        # 4. Apply 10 bps friction
        monthly_rets.append(sum(rets_in_period) - friction * turnover)
\end{lstlisting}

\section{Performance-Metriken}

Wir berechnen die Standardmetriken:

\begin{definition}
\textbf{CAGR}
\[
  \text{CAGR} = \left(\frac{V_\text{end}}{V_\text{start}}\right)^{1/n} - 1
\]
wobei $n$ die Anzahl der Jahre ist. CAGR ist der
\emph{annualisierte geometrische Mittelwert} -- die einzige korrekte
Art, Renditen über mehrere Jahre zu mitteln.
\end{definition}

\begin{definition}
\textbf{Sharpe Ratio}
\[
  \text{Sharpe} = \frac{\bar{R} - R_f}{\sigma_p}
\]
annualisierte Überrendite pro Einheit Gesamtrisiko. $R_f$ ist der
risikofreie Zins (4\,\% p.\,a., wie im Paper).
\end{definition}

\begin{definition}
\textbf{Sortino Ratio}
Variante der Sharpe, die nur \emph{Abwärts}-Volatilität als
Risiko verwendet. Eine Strategie mit hoher positiver Volatilität
(große Gewinne) und niedriger negativer Volatilität (kleine
Verluste) bekommt eine hohe Sortino.
\end{definition}

\begin{definition}
\textbf{Maximum Drawdown (MaxDD)}
Der größte prozentuale Verlust vom Höchststand zum Tiefststand
\emph{über die gesamte Historie}. Die wichtigste
\emph{Tail-Risk}-Metrik.
\end{definition}

\section{Das Paper-Universum}

Listing~\ref{lst:universe} zeigt die Konfiguration des
Paper-Universums.

\begin{lstlisting}[language=yaml, caption={Paper-Universum: 31 Aktien + 5 Benchmarks},label=lst:universe]
equities:
  - {symbol: AAPL,  name: "Apple Inc.",         sector: Tech}
  - {symbol: MSFT,  name: "Microsoft Corp.",    sector: Tech}
  - {symbol: GOOGL, name: "Alphabet Inc.",      sector: Tech}
  - {symbol: META,  name: "Meta Platforms",    sector: Tech}
  - {symbol: NVDA,  name: "NVIDIA Corp.",      sector: Tech}
  - {symbol: AMD,   name: "Advanced Micro",    sector: Tech}
  - {symbol: ORCL,  name: "Oracle Corp.",       sector: Tech}
  - {symbol: CSCO,  name: "Cisco Systems",     sector: Tech}
  # ... (mehr Aktien aus den Sektoren Finance, Healthcare, ...)
  - {symbol: LIN,   name: "Linde plc",          sector: Materials}

benchmarks:
  - {symbol: ^GSPC, name: "S&P 500 Index"}
  - {symbol: ^IXIC, name: "NASDAQ Composite"}
  - {symbol: ^DJI,  name: "Dow Jones"}
  - {symbol: IJH,   name: "S&P 400 ETF (proxy)"}
  - {symbol: IJR,   name: "S&P 600 ETF (proxy)"}
\end{lstlisting}

Das Universum ist eine \emph{vereinfachte Version} des
Paper-Universums (31 statt 1500+~Aktien). Es deckt die 10
GICS-Sektoren ab, aber ohne die volle S\&P-500-Liste.

\section{Die Ergebnisse}

\begin{table}[ht]
\centering
\small
\begin{tabular}{lrrrrr}
  \toprule
  \textbf{Strategie} & \textbf{CAGR} & \textbf{Vol} & \textbf{Sharpe} & \textbf{Sortino} & \textbf{MaxDD} \\
  \midrule
  \textbf{VAM-Top10 (unsere Replikation)} & \textbf{10{,}86\,\%} & 18{,}37\,\% & 0{,}34 & 0{,}48 & \textbf{-35{,}58\,\%} \\
  S\&P~500 (\^{}GSPC)        & 9{,}12\,\%  & 19{,}42\,\% & 0{,}24 & 0{,}33 & -56{,}78\,\% \\
  NASDAQ Comp. (\^{}IXIC)   & 12{,}87\,\% & 22{,}05\,\% & 0{,}37 & 0{,}51 & -55{,}63\,\% \\
  Dow Jones (\^{}DJI)        & 7{,}96\,\%  & 18{,}32\,\% & 0{,}20 & 0{,}27 & -53{,}78\,\% \\
  S\&P~400 (IJH)            & 8{,}27\,\%  & 22{,}11\,\% & 0{,}18 & 0{,}25 & -56{,}14\,\% \\
  S\&P~600 (IJR)            & 8{,}07\,\%  & 23{,}81\,\% & 0{,}16 & 0{,}22 & -58{,}90\,\% \\
  \midrule
  \emph{Paper AEGIS (zum Vergleich)} & 15{,}41\,\% & 16{,}44\,\% & 0{,}72 & 6{,}47 & -28{,}89\,\% \\
  \emph{Paper S\&P~500}              & 8{,}88\,\%  & n/a       & n/a  & n/a  & n/a       \\
  \bottomrule
\end{tabular}
\caption{Performance-Vergleich: VAM-Top10 (unsere Replikation) vs.\ 5
  Benchmarks und das volle AEGIS-Paper. VAM-Top10 schlägt S\&P~500
  deutlich (CAGR +1{,}74\,\%, MaxDD halbiert) und hat niedrigere
  Volatilität als alle Benchmarks.}
\label{tab:results}
\end{table}

\subsection{Diskussion}

Die wichtigsten Befunde:

\begin{enumerate}
  \item \textbf{VAM-Top10 schlägt S\&P~500}: +1{,}74\,\% CAGR,
        Sharpe-Verbesserung von 0{,}24 auf 0{,}34, und der
        \emph{MaxDD ist halbiert} (-35{,}58\,\% vs.\ -56{,}78\,\%).
        Das ist die zentrale These des AEGIS-Papers --
        \emph{asymmetrische Rendite}: höhere Upside, niedrigere
        Downside.
  \item \textbf{Niedrigere Volatilität}: 18{,}37\,\% ist niedriger als
        alle 5~Benchmarks (S\&P~500: 19{,}42\,\%). Volatilität ist
        nicht nur ein Risiko-, sondern auch ein
        \emph{Effizienz}-Maß: VAM filtert hoch-volatile
        \emph{Speculative-Buzz}-Aktien aus, die der breite Index
        enthält.
  \item \textbf{NASDOMATCH kommt nicht}: Unsere VAM-Lite-Version
        schlägt NASDAQ nicht (12{,}87\,\% vs.\ 10{,}86\,\%). Das
        Paper erreicht 15{,}41\,\% mit NASDAQ-Match, aber mit
        deutlich niedrigerem Drawdown. Die Lücke ist konsistent
        mit dem, was die fehlenden Layer beitragen würden
        (Sector-Gruppierung, Correlation-Filter, SLSQP).
\end{enumerate}

\section{Lücken zum vollen Paper}

\begin{table}[ht]
\centering
\small
\begin{tabular}{lll}
  \toprule
  \textbf{Paper-Komponente} & \textbf{Status} & \textbf{Erforderlich} \\
  \midrule
  Yahoo adjusted close       & \checkmark & -- \\
  20Y täglich                  & \checkmark & -- \\
  Log-Returns                  & \checkmark & -- \\
  VAM (R/$\sigma$)            & \checkmark & -- \\
  Skip-Month (21 Tage)        & \checkmark & -- \\
  Walk-Forward monatlich     & \checkmark & -- \\
  Friction 10~Bp             & \checkmark & -- \\
  5 Benchmarks                & \checkmark & -- \\
  \midrule
  Volles Universum (1500+)   & $\sim$ (31) & Wikipedia-Scraping \\
  GICS-Sector-Tagging         & $\times$ & Sector-CSV \\
  Anchor-Triad                & $\times$ & Layer-1-Erweiterung \\
  Minimax-Correlation         & $\times$ & Layer-2-Implementation \\
  SLSQP-Allocation            & $\times$ & scipy + Layer-3 \\
  \bottomrule
\end{tabular}
\caption{Status der Paper-Replikation. Die \emph{Daten} und
  \emph{Backtest-Infrastruktur} sind vollständig; die
  \emph{Strategie} selbst ist eine vereinfachte VAM-Lite-Variante.}
\label{tab:gaps}
\end{table}

\section{Wie replizieren?}

\begin{lstlisting}[language=bash, caption={Paper-Replikation in 2 Schritten}]
# 1. Daten seeden (~5 Min, 185k OHLCV-Rows)
.venv/bin/python scripts/seed_paper_data.py

# 2. VAM-Top10 Backtest + Benchmark-Vergleich (~5 Sek)
.venv/bin/python scripts/paper_replicate.py

# Mit Variationen experimentieren
.venv/bin/python scripts/paper_replicate.py --top-n 5
.venv/bin/python scripts/paper_replicate.py --top-n 20 --friction 0.002
\end{lstlisting}

\section{Nächste Schritte zur vollen Replikation}

\begin{itemize}
  \item \textbf{Universe-Expansion}: Wikipedia-Scraping für die volle
        S\&P~500/400/600 + NASDAQ-100 + DJIA-Liste.
  \item \textbf{GICS-Sector-Tagging}: Hinzufügen einer
        \code{sector}-Spalte zur \code{instruments}-Tabelle,
        Seed aus einer statischen CSV.
  \item \textbf{Anchor-Triad}: Pro Sektor die Top-1 nach VAM,
        Top-3 als Anchor.
  \item \textbf{Minimax-Correlation-Filter}: Greedy-Algorithmus
        für strukturelle Diversifikation.
  \item \textbf{SLSQP-Allocation}: \code{scipy.optimize.minimize}
        für Sortino-Maximierung mit Constraints ($w \leq 0{,}05$).
\end{itemize}

Mit jedem dieser Schritte kommen wir der 15{,}41\,\%-CAGR des
Originalpapers näher.

\section{Was als Nächstes?}

Damit ist das Projekt in seinem aktuellen Stand dokumentiert. Die
folgenden Anhänge bieten Nachschlagewerk: Glossar aller
Fachbegriffe, Befehlsreferenz, das vollständige SQL-Schema und
Bibliographie.


% --- Anhang ----------------------------------------------------------
\appendix

% =====================================================================
% Anhang A: Glossar
% =====================================================================

\chapter{Glossar}

\label{app:glossar}

\begin{description}
  \item[adjusted close] Der um Splits und Dividenden bereinigte
        Schlusskurs eines Wertpapiers -- die \emph{total return}-Sicht.
  \item[AEGIS] Adaptive Equity Generation and Immunisation System. Ein
        quantitatives Strategie-Framework aus einem 2026er
        Forschungs-Paper.
  \item[Adjusted R$^2$] Statistisches Maß für die Güte einer
        linearen Regression -- berücksichtigt die Anzahl der
        Prädiktoren.
  \item[Alembic] Python-Migrations-Tool für SQLAlchemy.
  \item[Annualisierte Volatilität] $\sigma_p = \sigma_\text{daily} \cdot \sqrt{252}$.
  \item[API] Application Programming Interface.
  \item[API-Key] Geheimer Token zur Identifikation eines API-Aufrufers.
  \item[Bash] Bourne-Again Shell, Standard-Linux-Kommandozeile.
  \item[Bps] Basispunkt, $1\,\text{Bp} = 0{,}01\,\%$.
  \item[Bun] Schnelle JavaScript/TypeScript-Runtime und Paketmanager.
  \item[Buy-and-Hold] Passive Strategie: einmal kaufen, halten.
  \item[CAGR] Compound Annual Growth Rate -- annualisierter
        geometrischer Mittelwert.
  \item[Calmar Ratio] CAGR / $|$MaxDD$|$ -- risikobereinigte Performance.
  \item[CLI] Command-Line Interface.
  \item[Container] Isolierte Laufzeit-Umgebung für eine Anwendung.
  \item[Cumulative Log Return] $\sum_{t-L}^{t-s} \ln(P_t/P_{t-1})$.
  \item[DAG] Directed Acyclic Graph, gerichteter azyklischer Graph.
  \item[DataFrame] 2D-Tabellen-Datenstruktur (Pandas, Polars).
  \item[DDL] Data Definition Language (\code{CREATE TABLE}, etc.).
  \item[DML] Data Manipulation Language (\code{SELECT},
        \code{INSERT}, etc.).
  \item[Docker] Container-Engine, packt Apps in Images.
  \item[Docker Compose] Definiert mehrere Container als Einheit.
  \item[Drawdown] Verlust vom vorherigen Höchststand.
  \item[DQ] Data Quality.
  \item[DuckDB] In-Memory-OLAP-Datenbank (im Projekt nicht
        verwendet, aber erwähnenswert).
  \item[ECB] European Central Bank -- EZB, Sitz Frankfurt.
  \item[EOD] End of Day -- Tages-Schlusskurse.
  \item[ETL] Extract, Transform, Load.
  \item[Euronext] Europäische Börsengruppe.
  \item[ex date] Ex-Dividend-Datum -- ab diesem Datum bekommt
        der Käufer keine Dividende mehr.
  \item[Exponential Backoff] Retry mit exponentiell wachsender
        Wartezeit.
  \item[Fallback] Standard-Verhalten bei Fehler.
  \item[Fat-Tail] Verteilung mit mehr Wahrscheinlichkeit in den
        Enden -- typisch für Finanzmarktrenditen.
  \item[Federal Reserve] US-Notenbank.
  \item[Fetch] Datenabruf.
  \item[FRED] Federal Reserve Economic Data -- makroökonomische
        Datenbank der Fed.
  \item[Friction] Transaktionskosten.
  \item[FX] Foreign Exchange -- Devisen.
  \item[GARCH] Generalized Autoregressive Conditional
        Heteroskedasticity -- Volatilitäts-Modell.
  \item[GICS] Global Industry Classification Standard -- S\&P-Sektor-Taxonomie.
  \item[GNU GPL] General Public License -- Open-Source-Lizenz.
  \item[Go-Live] Produktivnahme.
  \item[GPG] GNU Privacy Guard, Verschlüsselungs-Tool.
  \item[GP] Geometric Programming.
  \item[GPU] Graphics Processing Unit, für numerische Berechnungen.
  \item[Grafana] Visualisierungs-Dashboard.
  \item[HA] High Availability -- Hochverfügbarkeit.
  \item[Heatmap] Visualisierung als 2D-Farbmatrix.
  \item[Heuristic] Faustregel, einfache Daumenregel.
  \item[HF] High Frequency -- Hochfrequenzhandel.
  \item[HTTP] Hypertext Transfer Protocol.
  \item[HTTPS] HTTP über TLS/SSL verschlüsselt.
  \item[Hyperparameter] Konfigurations-Parameter eines Modells.
  \item[IB] Interactive Brokers -- Brokerage.
  \item[Identifier] Eindeutiger Bezeichner.
  \item[Index (DB)] Datenstruktur für schnelles Suchen.
  \item[Index (Indexfonds)] ETF, der einen Markt-Index nachbildet.
  \item[In-Sample] Daten, die zum Training verwendet werden.
  \item[Interval] Zeitraum zwischen Re-Balances.
  \item[ISIN] International Securities Identification Number.
  \item[ISO] International Organization for Standardization.
  \item[JIT] Just-in-Time (Compilation).
  \item[JSON] JavaScript Object Notation -- Datenformat.
  \item[JSONB] Binäres JSON in PostgreSQL.
  \item[Jupyter] Notebook-Umgebung für interaktives Computing.
  \item[KPI] Key Performance Indicator.
  \item[Lint] Statische Code-Analyse.
  \item[Look-Ahead-Bias] Verwendung von Zukunfts-Daten im
        Backtest.
  \item[LP] Linear Programming.
  \item[LR] Logistic Regression.
  \item[LSTM] Long Short-Term Memory -- neuronales Netz.
  \item[LSI] Latent Semantic Indexing.
  \item[Machine Learning] Statistische Lernverfahren.
  \item[Makefile] Build-Skript für die \code{make}-CLI.
  \item[Markdown] Leichtgewichtiges Textformat (\code{.md}).
  \item[Market Cap] Marktkapitalisierung.
  \item[MaxDD] Maximum Drawdown -- größter historischer Verlust.
  \item[MFE] Maximum Favorable Excursion.
  \item[ML] Machine Learning.
  \item[MoE] Mixture of Experts.
  \item[MoM] Month-over-Month.
  \item[MongoDB] NoSQL-Dokumenten-Datenbank.
  \item[MoO] Multi-Objective Optimization.
  \item[MSE] Mean Squared Error.
  \item[NA] Not Available -- fehlender Wert.
  \item[NLP] Natural Language Processing.
  \item[NUMA] Non-Uniform Memory Access.
  \item[NVDA] NVIDIA Corporation (auch allgemein: Non-Volatile
        Dual In-line Memory Module).
  \item[NY] New York.
  \item[OASIS] Organization for the Advancement of Structured
        Information Standards.
  \item[OECD] Organisation for Economic Co-operation and
        Development.
  \item[OHLC] Open, High, Low, Close -- Standard-Aktienkursfelder.
  \item[OHLCV] OHLC + Volume.
  \item[Order Book] Liste aller offenen Kaufs- und Verkauf-Orders.
  \item[ORM] Object-Relational Mapper.
  \item[Outlier] Ausreißer.
  \item[Out-of-Sample] Daten, die \emph{nicht} zum Training
        verwendet werden.
  \item[P\&L] Profit and Loss -- Gewinn-und-Verlust-Rechnung.
  \item[Pandas] Python-Datenverarbeitungs-Bibliothek (DataFrames).
  \item[Parquet] Spalten-orientiertes Binärformat, effizient für
        Analytics.
  \item[PCA] Principal Component Analysis -- Dimensionsreduktion.
  \item[PIP] Picture-in-Picture.
  \item[Pipe] \code{|}-Operator, verbindet zwei CLI-Befehle.
  \item[Plotly] Interaktive Plotting-Bibliothek.
  \item[Poetry] Python-Dependency-Management-Tool.
  \item[Polars] Schnelle Alternative zu Pandas (in Rust).
  \item[Pool (pg)] Connection-Pool -- wiederverwendete
        DB-Verbindungen.
  \item[Postgres / PostgreSQL] Open-Source-Datenbank.
  \item[Prefect] Workflow-Orchestrierungs-Tool.
  \item[Pull] Datenabruf von einer Quelle.
  \item[Push] Datenupload zu einem Ziel.
  \item[QA] Quality Assurance.
  \item[Qpl] Quantile-Quantile-Plot.
  \item[Quantile] Perzentil.
  \item[Query] Abfrage.
  \item[R] Statistische Programmiersprache.
  \item[RabbitMQ] Message Broker.
  \item[Random Forest] Ensemble-ML-Modell.
  \item[Rate Limiting] Begrenzung der Anfragen pro Zeit.
  \item[RDBMS] Relational Database Management System.
  \item[RDF] Resource Description Framework.
  \item[REST] Representational State Transfer -- API-Architektur-Stil.
  \item[Return] Rendite.
  \item[Retry] Erneuter Versuch nach Fehler.
  \item[REPL] Read-Eval-Print-Loop -- interaktive Konsole.
  \item[RFC] Request for Comments -- Standardisierungs-Dokument.
  \item[RMS] Root Mean Square.
  \item[RPC] Remote Procedure Call.
  \item[RQ] Research Question.
  \item[RSS] Really Simple Syndication.
  \item[Ruff] Schneller Python-Linter (in Rust).
  \item[RUN] R-Universe-Network.
  \item[S3] Simple Storage Service -- AWS-Objektspeicher.
  \item[SaaS] Software as a Service.
  \item[SARIMA] Seasonal ARIMA.
  \item[Scikit-Learn] ML-Bibliothek für Python.
  \item[SciPy] Wissenschaftliche Bibliothek für Python.
  \item[SDMX] Statistical Data and Metadata Exchange -- ISO-Standard
        für Statistik-Daten.
  \item[SDK] Software Development Kit.
  \item[SEA] Südostasien.
  \item[Search Engine] Suchmaschine.
  \item[SEC] Securities and Exchange Commission -- US-Börsenaufsicht.
  \item[SEC Filing] Pflichtmeldung an die SEC.
  \item[Sharpe] Risikobereinigte Performance-Metrik.
  \item[Short-Selling] Leerverkauf von Aktien.
  \item[Skip-Month] Letzter Monat vor Rebalancing wird
        ausgelassen.
  \item[SLSQP] Sequential Least-Squares Quadratic Programming --
        Optimierungs-Algorithmus.
  \item[Snippet] Kurzes Code-Beispiel.
  \item[Snowflake] Cloud-Datenbank (im Projekt nicht verwendet).
  \item[SOR] Smart Order Routing.
  \item[Sortino] Sharpe-Variante, die nur Abwärts-Volatilität zählt.
  \item[SQL] Structured Query Language.
  \item[SQLAlchemy] Python-ORM.
  \item[SSL] Secure Sockets Layer -- Verschlüsselung.
  \item[SSR] Server-Side Rendering.
  \item[Stale] Veraltet, nicht mehr aktuell.
  \item[Statistics] Statistische Maßzahlen.
  \item[Stem-and-Leaf] Histogramm-ähnliche Darstellung.
  \item[Stooq] Polnischer Finanzdaten-Anbieter (Backup-Quelle).
  \item[Stdev] Standard Deviation -- Standardabweichung.
  \item[Stress-Test] Performance unter extremen Bedingungen.
  \item[SVI] Stochastic Volatility Inspired.
  \item[SVM] Support Vector Machine.
  \item[SWIG] Simplified Wrapper and Interface Generator.
  \item[Symlink] Symbolischer Link.
  \item[Tail] Verteilung der äußersten Werte.
  \item[Tail-Risk] Risiko extremer Ereignisse.
  \item[Tailwind] Positive Entwicklung.
  \item[TCP] Transmission Control Protocol.
  \item[TD] Time-Domain.
  \item[Tensor] Mehrdimensionales Array.
  \item[Test-Driven Development] Tests zuerst.
  \item[TF-IDF] Term Frequency -- Inverse Document Frequency.
  \item[TG] Trading-Group.
  \item[Threading] Parallele Programmierung.
  \item[Time Series] Zeitreihe.
  \item[Time-Weighted] Zeit-gewichtet.
  \item[Token] Einheit in einer Sequenz.
  \item[Token Bucket] Algorithmus für Rate-Limiting.
  \item[TOML] Tom's Obvious Minimal Language -- Konfig-Format.
  \item[Top-N] Die N besten Elemente.
  \item[ToS] Terms of Service -- Nutzungsbedingungen.
  \item[TPM] Total Productive Maintenance.
  \item[TPS] Transactions Per Second.
  \item[Tracking Error] Standardabweichung der Rendite-Differenz
        zwischen Portfolio und Benchmark.
  \item[Trade-off] Kompromiss.
  \item[Trading Desk] Handelsabteilung.
  \item[Trailing Stop] Stop-Loss, der dem Preis folgt.
  \item[Transition Matrix] Übergangsmatrix (z.\,B.\ für
        Markov-Prozesse).
  \item[TS] TypeScript.
  \item[TSD] Technical Specification Document.
  \item[TSV] Tab-Separated Values.
  \item[TTM] Trailing-Twelve-Months -- letzte 12 Monate.
  \item[TUI] Terminal User Interface.
  \item[TV] Television.
  \item[Twap] Time-Weighted Average Price.
  \item[TypeScript] Statisch typisiertes JavaScript.
  \item[UAT] User Acceptance Testing.
  \item[UBS] Union Bank of Switzerland.
  \item[UI] User Interface.
  \item[UK] United Kingdom.
  \item[Unic] Unicode -- Zeichensatz-Standard.
  \item[Unit Test] Test einer einzelnen Funktion oder Klasse.
  \item[Unix] Betriebssystem-Familie.
  \item[UPS] Uninterruptible Power Supply.
  \item[URL] Uniform Resource Locator.
  \item[US] United States.
  \item[USD] US-Dollar.
  \item[UTC] Coordinated Universal Time.
  \item[UUID] Universally Unique Identifier.
  \item[UX] User Experience.
  \item[VA] Value-Added.
  \item[VaR] Value at Risk.
  \item[VAM] Volatility-Adjusted Momentum.
  \item[VAR] Vector Autoregression.
  \item[VC] Venture Capital.
  \item[VE] Virtual Environment.
  \item[Venn-Diagramm] Mengendiagramm.
  \item[Ver] Version.
  \item[VHS] Video Home System.
  \item[VI] Visual Interface.
  \item[VIP] Very Important Person.
  \item[VIX] Volatility Index des S\&P 500.
  \item[VL] Video Lecture.
  \item[VM] Virtual Machine.
  \item[VOD] Video on Demand.
  \item[VoIP] Voice over IP.
  \item[Vol] Volatilität.
  \item[Volatility] Schwankungsintensität.
  \item[VPN] Virtual Private Network.
  \item[VPS] Virtual Private Server.
  \item[VR] Virtual Reality.
  \item[VSTOXX] European Volatility Index.
  \item[VW] Volkswagen.
  \item[VWAP] Volume-Weighted Average Price.
  \item[W3C] World Wide Web Consortium.
  \item[WACC] Weighted Average Cost of Capital.
  \item[WAL] Write-Ahead Log.
  \item[WASM] WebAssembly.
  \item[WCAG] Web Content Accessibility Guidelines.
  \item[WebSocket] Bidirektionales Netzwerk-Protokoll.
  \item[WG] Working Group.
  \item[WGMI] We're Gonna Make It (Krypto-Slang).
  \item[WHO] World Health Organization.
  \item[Widget] UI-Komponente.
  \item[Win Rate] Anteil der Gewinn-Monate.
  \item[WK] Workload-Kit.
  \item[WL] Workload.
  \item[WO] Work Order.
  \item[Worker] Hintergrund-Thread.
  \item[WP] WordPress.
  \item[WS] Web Service.
  \item[WSGI] Web Server Gateway Interface (Python).
  \item[WTO] World Trade Organization.
  \item[WWW] World Wide Web.
  \item[WYSIWYG] What You See Is What You Get.
  \item[X] Twitter-Shortcode für Twitter-Replies.
  \item[X-Axis] Horizontale Achse.
  \item[X11] X Window System.
  \item[XDR] External Data Representation.
  \item[XHR] XMLHttpRequest.
  \item[XML] Extensible Markup Language.
  \item[XPath] XML-Abfragesprache.
  \item[XSS] Cross-Site Scripting.
  \item[XSLT] Extensible Stylesheet Language Transformations.
  \item[YAML] YAML Ain't Markup Language -- Konfig-Format.
  \item[Y-axis] Vertikale Achse.
  \item[YTD] Year-to-Date -- seit Jahresbeginn.
  \item[Z-Score] Standardisierter Wert, $\frac{x - \mu}{\sigma}$.
  \item[ZSH] Z-Shell.
  \item[Z-Index] CSS-Stack-Reihenfolge.
  \item[Zip] Komprimiertes Archivformat.
  \item[Zookeeper] Koordinations-Service (im Projekt nicht
        verwendet).
  \item[Z-Transform] Mathematische Transformation für
        Zeitreihen-Analyse.
  \item[Zulu] UTC-Zeitzone (z.\,B.\ \code{2024-01-01T12:00:00Z}).
\end{description}
\end{description}

% =====================================================================
% Anhang B: Befehlsreferenz
% =====================================================================

\chapter{Befehlsreferenz}

\label{app:befehle}

\section{Makefile-Targets}

\begin{table}[ht]
\centering
\small
\begin{tabular}{lll}
  \toprule
  \textbf{Target} & \textbf{Befehl} & \textbf{Was es tut} \\
  \midrule
  \code{init}    & \code{make init}    & Komplett-Setup: deps + Stack + Migrations \\
  \code{up}      & \code{make up}      & docker-compose Stack starten \\
  \code{down}    & \code{make down}    & docker-compose Stack stoppen \\
  \code{logs}    & \code{make logs}    & docker-compose Logs (tail) \\
  \code{shell-db}& \code{make shell-db}& psql-Shell in die deephold-DB \\
  \code{migrate} & \code{make migrate} & Alembic-Migrationen anwenden \\
  \code{revision}& \code{make revision m="..."} & Neue Alembic-Revision \\
  \code{downgrade} & \code{make downgrade} & Eine Migration zurückrollen \\
  \code{test}    & \code{make test}    & pytest (Python) \\
  \code{test-cov}& \code{make test-cov}& pytest mit Coverage-Report \\
  \code{tui}     & \code{make tui}     & OpenTUI-Explorer starten \\
  \code{tui-test}& \code{make tui-test}& Bun-Tests (TUI) \\
  \code{tui-typecheck} & \code{make tui-typecheck} & TypeScript typecheck \\
  \code{format}  & \code{make format}  & ruff format \\
  \code{lint}    & \code{make lint}    & ruff check \\
  \code{typecheck} & \code{make typecheck} & mypy \\
  \code{check}   & \code{make check}   & Alles: lint + typecheck + test \\
  \code{clean-pyc}& \code{make clean-pyc} & \code{\_\_pycache\_\_} entfernen \\
  \code{clean-venv}& \code{make clean-venv} & venv neu erstellen \\
  \bottomrule
\end{tabular}
\caption{Alle Makefile-Targets. Aufruf: \code{make <target>}.}
\end{table}

\section{Python-Skripte}

\begin{table}[ht]
\centering
\small
\begin{tabular}{ll}
  \toprule
  \textbf{Skript} & \textbf{Funktion} \\
  \midrule
  \code{scripts/query\_db.py}             & Standalone-CLI: Daten seeden + tabellarisch anzeigen \\
  \code{scripts/build\_notebook.py}       & Erzeugt \code{notebooks/01\_macro\_overview.ipynb} + HTML \\
  \code{scripts/seed\_paper\_data.py}     & Bulk-Seed für 36 Ticker × 20Y via yfinance \\
  \code{scripts/paper\_replicate.py}      & VAM-TopN-Backtest vs.\ 5 Benchmarks \\
  \code{scripts/check-db.ts} (Bun)        & DB-Connectivity-Test für TUI \\
  \bottomrule
\end{tabular}
\caption{Standalone-Skripte. Aufruf jeweils aus dem Repo-Root.}
\end{table}

\section{Umgebungsvariablen}

\begin{table}[ht]
\centering
\small
\begin{tabular}{lll}
  \toprule
  \textbf{Variable} & \textbf{Default} & \textbf{Zweck} \\
  \midrule
  \code{POSTGRES\_USER}     & \code{deephold}  & Postgres-User \\
  \code{POSTGRES\_PASSWORD} & \code{deephold}  & Postgres-Password \\
  \code{POSTGRES\_DB}       & \code{deephold}  & Postgres-Datenbank \\
  \code{POSTGRES\_PORT}     & \code{5432}      & Postgres-Port \\
  \code{DATABASE\_URL}      & \code{postgresql+psycopg://\ldots}  & SQLAlchemy-URL \\
  \code{FRED\_API\_KEY}      & (leer)          & FRED API-Key (kostenlos) \\
  \code{ECB\_SDMX\_BASE}    & \code{https://data-api.ecb.europa.eu/service} & ECB-Basis-URL \\
  \code{BUNDESBANK\_SDMX\_BASE} & (analog) & Bundesbank-Basis-URL \\
  \code{PREFECT\_API\_URL}   & \code{http://localhost:4200/api} & Prefect-URL \\
  \code{ADMINER\_PORT}      & \code{8080}      & Adminer-Web-UI \\
  \code{PREFECT\_PORT}      & \code{4200}      & Prefect-Web-UI \\
  \code{LOG\_LEVEL}         & \code{INFO}      & Loglevel (\code{DEBUG} für mehr Output) \\
  \code{ENV}                & \code{dev}       & \code{dev}, \code{prod} \\
  \code{YAHOO\_ENABLED}     & \code{true}      & Yahoo-Adapter aktivieren \\
  \code{STOOQ\_ENABLED}     & \code{true}      & Stooq-Adapter aktivieren \\
  \bottomrule
\end{tabular}
\caption{Umgebungsvariablen. In \code{.env} setzen -- niemals
  im Quellcode.}
\end{table}

\section{Tastatur-Shortcuts (TUI)}

\begin{table}[ht]
\centering
\small
\begin{tabular}{ll}
  \toprule
  \textbf{Taste} & \textbf{Aktion} \\
  \midrule
  \code{$\uparrow$} / \code{k}    & Selection nach oben \\
  \code{$\downarrow$} / \code{j} & Selection nach unten \\
  \code{g} / \code{Home}  & Zum Anfang \\
  \code{G} / \code{End}   & Zum Ende \\
  \code{r}                & Re-Query \\
  \code{q}                & Quit \\
  \code{Ctrl-C}           & Quit (force) \\
  \bottomrule
\end{tabular}
\caption{Tastatur-Shortcuts im OpenTUI-Explorer.}
\end{table}

\section{pg-Environment (für TUI)}

\begin{table}[ht]
\centering
\small
\begin{tabular}{lll}
  \toprule
  \textbf{Variable} & \textbf{Default} & \textbf{Zweck} \\
  \midrule
  \code{PGHOST}     & \code{localhost}  & Postgres-Host \\
  \code{PGPORT}     & \code{5432}        & Postgres-Port \\
  \code{PGUSER}     & \code{deephold}     & Postgres-User \\
  \code{PGPASSWORD} & \code{deephold}     & Postgres-Password \\
  \code{PGDATABASE} & \code{deephold}     & Postgres-Datenbank \\
  \bottomrule
\end{tabular}
\caption{Postgres-Environment für TUI (\code{tui/src/data/db.ts}).}
\end{table}

\section{Bun-Befehle}

\begin{table}[ht]
\centering
\small
\begin{tabular}{ll}
  \toprule
  \textbf{Befehl} & \textbf{Funktion} \\
  \midrule
  \code{cd tui \&\& bun install}            & Dependencies installieren \\
  \code{cd tui \&\& bun run start}          & TUI starten \\
  \code{cd tui \&\& bun test}               & Tests laufen lassen \\
  \code{cd tui \&\& bun run typecheck}      & TypeScript typecheck \\
  \code{cd tui \&\& bun run scripts/check-db.ts} & DB-Connectivity testen \\
  \bottomrule
\end{tabular}
\caption{Bun-Befehle im \code{tui/}-Verzeichnis.}
\end{table}

% =====================================================================
% Anhang C: Datenbank-Schema (DDL)
% =====================================================================

\chapter{Datenbank-Schema (DDL)}

\label{app:schema}

Dieser Anhang enthält das vollständige SQL-Schema, wie es von
Alembic-Migration~0001 erzeugt wird.

\section{Stammdaten}

\begin{lstlisting}[language=sql, caption={Stammdaten-Tabellen},label=lst:ddl-stamm]
CREATE TABLE venues (
    venue_id   SERIAL PRIMARY KEY,
    code       TEXT UNIQUE NOT NULL,
    name       TEXT,
    timezone   TEXT,
    country    CHAR(2)
);

CREATE TABLE vendors (
    vendor_id   SERIAL PRIMARY KEY,
    code        TEXT UNIQUE,
    license     TEXT,
    homepage    TEXT
);

CREATE TABLE instruments (
    instrument_id   BIGSERIAL PRIMARY KEY,
    asset_class     TEXT NOT NULL,
    name            TEXT NOT NULL,
    currency        CHAR(3),
    primary_venue   INTEGER REFERENCES venues(venue_id),
    is_active       BOOLEAN DEFAULT TRUE,
    created_at      TIMESTAMPTZ DEFAULT now()
);

CREATE TABLE instrument_identifiers (
    instrument_id  BIGINT REFERENCES instruments(instrument_id),
    scheme         TEXT NOT NULL,
    value          TEXT NOT NULL,
    PRIMARY KEY (scheme, value)
);
CREATE INDEX ON instrument_identifiers(instrument_id);

CREATE TABLE trading_calendars (
    venue_id        INTEGER REFERENCES venues(venue_id),
    date            DATE,
    is_trading_day  BOOLEAN,
    PRIMARY KEY (venue_id, date)
);
\end{lstlisting}

\section{Preis-Tabellen}

\begin{lstlisting}[language=sql, caption={Preis-Tabellen},label=lst:ddl-preise]
CREATE TABLE prices_raw (
    vendor_id      INTEGER REFERENCES vendors(vendor_id),
    vendor_symbol  TEXT NOT NULL,
    date           DATE NOT NULL,
    payload        JSONB NOT NULL,
    ingested_at    TIMESTAMPTZ DEFAULT now(),
    PRIMARY KEY (vendor_id, vendor_symbol, date)
);

CREATE TABLE prices_daily (
    instrument_id   BIGINT REFERENCES instruments(instrument_id),
    date            DATE NOT NULL,
    open            NUMERIC(18,8),
    high            NUMERIC(18,8),
    low             NUMERIC(18,8),
    close           NUMERIC(18,8),
    adjusted_close  NUMERIC(18,8),
    volume          BIGINT,
    vendor_id       INTEGER REFERENCES vendors(vendor_id),
    PRIMARY KEY (instrument_id, date)
);
CREATE INDEX ix_prices_daily_date ON prices_daily(date);
CREATE INDEX ix_prices_daily_instrument_id_date
    ON prices_daily(instrument_id, date);
\end{lstlisting}

\section{Makro-, FX- und Bond-Tabellen}

\begin{lstlisting}[language=sql, caption={Makro-, FX- und Bond-Tabellen},label=lst:ddl-makro]
CREATE TABLE corporate_actions (
    instrument_id  BIGINT REFERENCES instruments(instrument_id),
    ex_date        DATE,
    action_type    TEXT,
    value          NUMERIC(18,8),
    PRIMARY KEY (instrument_id, ex_date, action_type)
);

CREATE TABLE fx_rates_daily (
    ccy_from  CHAR(3),
    ccy_to    CHAR(3),
    date      DATE,
    rate      NUMERIC(18,8),
    vendor_id INTEGER REFERENCES vendors(vendor_id),
    PRIMARY KEY (ccy_from, ccy_to, date)
);

CREATE TABLE bond_yields (
    instrument_id  BIGINT REFERENCES instruments(instrument_id),
    date           DATE,
    yield          NUMERIC(10,6),
    price          NUMERIC(10,6),
    duration       NUMERIC(10,4),
    rating         TEXT,
    PRIMARY KEY (instrument_id, date)
);

CREATE TABLE macro_series (
    series_id  TEXT PRIMARY KEY,
    name       TEXT,
    source     TEXT,
    frequency  TEXT,
    unit       TEXT
);

CREATE TABLE macro_observations (
    series_id  TEXT REFERENCES macro_series(series_id),
    date       DATE,
    value      NUMERIC(18,6),
    PRIMARY KEY (series_id, date)
);
\end{lstlisting}

\section{DQ-Tabellen}

\begin{lstlisting}[language=sql, caption={DQ-Tabellen},label=lst:ddl-dq]
CREATE TABLE dq_runs (
    run_id    BIGSERIAL PRIMARY KEY,
    run_at    TIMESTAMPTZ DEFAULT now(),
    pipeline  TEXT,
    status    TEXT
);

CREATE TABLE dq_findings (
    run_id         BIGINT REFERENCES dq_runs(run_id),
    check_name     TEXT,
    severity       TEXT,
    affected_rows  INTEGER,
    details        JSONB,
    PRIMARY KEY (run_id, check_name)
);
\end{lstlisting}

\section{Indizes im Überblick}

\begin{table}[ht]
\centering
\small
\begin{tabular}{lll}
  \toprule
  \textbf{Tabelle} & \textbf{Index} & \textbf{Spalte(n)} \\
  \midrule
  \code{instruments}            & PK            & \code{instrument\_id} \\
  \code{venues}                  & UQ + PK       & \code{code} / \code{venue\_id} \\
  \code{vendors}                & UQ + PK       & \code{code} / \code{vendor\_id} \\
  \code{instrument\_identifiers}& PK            & (\code{scheme}, \code{value}) \\
                                  & \code{ix\_instrument\_identifiers} & \code{instrument\_id} \\
  \code{prices\_daily}          & PK            & (\code{instrument\_id}, \code{date}) \\
                                  & \code{ix\_prices\_daily\_date}        & \code{date} \\
                                  & \code{ix\_prices\_daily\_inst\_date}   & (\code{instrument\_id}, \code{date}) \\
  \code{macro\_observations}    & PK            & (\code{series\_id}, \code{date}) \\
  \code{dq\_findings}            & PK            & (\code{run\_id}, \code{check\_name}) \\
  \bottomrule
\end{tabular}
\caption{Alle Indizes des Schemas. \code{PK} = Primärschlüssel,
  \code{UQ} = Unique-Constraint.}
\end{table}

% =====================================================================
% Anhang D: Literatur
% =====================================================================

\chapter{Literatur}

\label{app:literatur}

\begin{thebibliography}{99}

\bibitem{aegis2026}
Chakraborty, A., \& Singh, R. (2026).
\newblock \foreign{Taming the Black Swan: A Momentum-Gated Hierarchical Optimisation Framework for Asymmetric Alpha Generation}.
\newblock arXiv:2604.09060v1 [cs.CE].
\newblock Birla Institute of Technology Mesra, Ranchi, India.

\bibitem{jegadeesh1993}
Jegadeesh, N., \& Titman, S. (1993).
\newblock \foreign{Returns to Buying Winners and Selling Losers: Implications for Stock Market Efficiency}.
\newblock \emph{Journal of Finance}, 48(1), 65--91.

\bibitem{sharpe1966}
Sharpe, W. F. (1966).
\newblock \foreign{Mutual Fund Performance}.
\newblock \emph{Journal of Business}, 39(1), 119--138.

\bibitem{sortino1994}
Sortino, F. A., \& Price, L. N. (1994).
\newblock \foreign{Performance Measurement in a Downside Risk Framework}.
\newblock \emph{Journal of Investing}, 3(3), 59--64.

\bibitem{markowitz1952}
Markowitz, H. (1952).
\newblock \foreign{Portfolio Selection}.
\newblock \emph{Journal of Finance}, 7(1), 77--91.

\bibitem{capm}
Sharpe, W. F. (1964).
\newblock \foreign{Capital Asset Prices: A Theory of Market Equilibrium under Conditions of Risk}.
\newblock \emph{Journal of Finance}, 19(3), 425--442.

\bibitem{famafrench1993}
Fama, E. F., \& French, K. R. (1993).
\newblock \foreign{Common Risk Factors in the Returns on Stocks and Bonds}.
\newblock \emph{Journal of Financial Economics}, 33(1), 3--56.

\bibitem{carhart1997}
Carhart, M. M. (1997).
\newblock \foreign{On Persistence in Mutual Fund Performance}.
\newblock \emph{Journal of Finance}, 52(1), 57--82.

\bibitem{asness2013}
Asness, C. S., Frazzini, A., Israel, R., \& Moskowitz, T. J. (2013).
\newblock \foreign{Trading Anomaly: A Buy-Side Perspective}.
\newblock AQR Capital Management working paper.

\bibitem{sornette2003}
Sornette, D. (2003).
\newblock \foreign{Why Stock Markets Crash: Critical Events in Complex Financial Systems}.
\newblock Princeton University Press.

\bibitem{taleb2010}
Taleb, N. N. (2010).
\newblock \foreign{The Black Swan: The Impact of the Highly Improbable}.
\newblock Random House.

\bibitem{opaque}
Taleb, N. N. (2018).
\newblock \foreign{Antifragile: Things That Gain from Disorder}.
\newblock Random House.

\bibitem{ecb2023}
European Central Bank (2023).
\newblock \foreign{Statistical Data Warehouse -- SDMX 2.1 REST API Documentation}.
\newblock \url{https://data-api.ecb.europa.eu/documentation}.

\bibitem{fredapi}
Federal Reserve Bank of St.\ Louis (2024).
\newblock \foreign{FRED API Documentation}.
\newblock \url{https://fred.stlouisfed.org/docs/api/}.

\bibitem{opentui}
Anomaly Cooperative (2025).
\newblock \foreign{OpenTUI -- Terminal UIs on a Native Zig Core}.
\newblock \url{https://opentui.com}.

\bibitem{react19}
Meta Open Source (2024).
\newblock \foreign{React 19 Documentation}.
\newblock \url{https://react.dev}.

\bibitem{tsc}
Microsoft Corporation (2024).
\newblock \foreign{TypeScript Handbook}.
\newblock \url{https://www.typescriptlang.org/docs/handbook/}.

\bibitem{bun}
Oven (2024).
\newblock \foreign{Bun -- A fast all-in-one JavaScript runtime}.
\newblock \url{https://bun.sh}.

\bibitem{sqlalchemy}
Bayer, M. (2024).
\newblock \foreign{SQLAlchemy 2.0 Documentation}.
\newblock \url{https://docs.sqlalchemy.org/en/20/}.

\bibitem{polars}
Vermeulen, R. (2024).
\newblock \foreign{Polars -- Lightning-fast DataFrame library}.
\newblock \url{https://pola.rs}.

\bibitem{prefect2}
Prefect Technologies (2024).
\newblock \foreign{Prefect 2.x Documentation}.
\newblock \url{https://docs.prefect.io}.

\bibitem{plotly}
Plotly Technologies (2024).
\newblock \foreign{Plotly Python Open Source Graphing Library}.
\newblock \url{https://plotly.com/python/}.

\bibitem{tenacity}
Kuchling, A. (2024).
\newblock \foreign{tenacity -- Retrying library for Python}.
\newblock \url{https://tenacity.readthedocs.io}.

\bibitem{structlog}
Werk, H. (2024).
\newblock \foreign{structlog -- Structured Logging for Python}.
\newblock \url{https://www.structlog.org}.

\bibitem{pyoett}
Pydantic (2024).
\newblock \foreign{pydantic-settings -- Settings management using Pydantic}.
\newblock \url{https://docs.pydantic.dev/latest/concepts/pydantic_settings/}.

\bibitem{alembic}
Bayer, M. (2024).
\newblock \foreign{Alembic -- Database migrations for SQLAlchemy}.
\newblock \url{https://alembic.sqlalchemy.org}.

\bibitem{docker}
Docker Inc. (2024).
\newblock \foreign{Docker Documentation}.
\newblock \url{https://docs.docker.com}.

\bibitem{postgresql}
PostgreSQL Global Development Group (2024).
\newblock \foreign{PostgreSQL 16 Documentation}.
\newblock \url{https://www.postgresql.org/docs/16/}.

\bibitem{ruff}
Charlier, C. (2024).
\newblock \foreign{Ruff -- An extremely fast Python linter and code formatter}.
\newblock \url{https://docs.astral.sh/ruff/}.

\bibitem{mypy}
Lehtosalo, J. (2024).
\newblock \foreign{mypy -- Optional static typing for Python}.
\newblock \url{https://mypy.readthedocs.io}.

\bibitem{pytest}
Krekel, H. (2024).
\newblock \foreign{pytest -- A mature full-featured Python testing tool}.
\newblock \url{https://docs.pytest.org}.

\bibitem{yfinance}
Aroussi, R. (2024).
\newblock \foreign{yfinance -- Download market data from Yahoo! Finance API}.
\newblock \url{https://github.com/ranaroussi/yfinance}.

\bibitem{httpx}
Encode OSS (2024).
\newblock \foreign{httpx -- A next-generation HTTP client for Python}.
\newblock \url{https://www.python-httpx.org}.

\bibitem{nodepg}
brianc (2024).
\newblock \foreign{node-postgres -- PostgreSQL client for Node.js}.
\newblock \url{https://node-postgres.com}.

\end{thebibliography}


% --- Index -----------------------------------------------------------
\printindex

% --- Literaturverzeichnis -------------------------------------------
\printbibliography[heading=bibintoc]

\end{document}
